\documentclass[10pt,a4paper]{article}
% ---------- packages ----------
\usepackage[margin=1.5cm]{geometry}
\usepackage{amsmath,amssymb,amsthm}
\usepackage{mathtools}
\usepackage{hyperref}
\hypersetup{
  colorlinks=true,
  linkcolor=black,          % internal refs stay black -- keeps the page calm
  citecolor=black,
  urlcolor=[rgb]{0,0.25,0.55},  % external URLs: visible, sober, printable
  breaklinks=true,
  pdftitle={Cross-Geometry Generalization of DoMINO for External Aerodynamics},
  pdfsubject={Working Project Plan}
}
\urlstyle{same}             % URLs in body font, not monospace
\usepackage{enumitem}
\usepackage{booktabs}
\usepackage{array}
\usepackage{fancyhdr}
\usepackage{titlesec}
\usepackage{microtype}
\usepackage{listings}
\usepackage{cleveref}
\usepackage{xcolor}

% ---------- compact typography ----------
\linespread{1.02}
\setlist{nosep,leftmargin=1.4em,topsep=2pt}
\setlength{\parskip}{2pt}
\setlength{\parindent}{0pt}

% ---------- section headings ----------
\titleformat{\section}{\normalsize\bfseries}{\thesection}{0.6em}{}
\titleformat{\subsection}{\small\bfseries}{\thesubsection}{0.6em}{}
\titleformat{\subsubsection}{\small\itshape}{\thesubsubsection}{0.6em}{}
\titlespacing*{\section}{0pt}{8pt}{3pt}
\titlespacing*{\subsection}{0pt}{6pt}{2pt}
\titlespacing*{\subsubsection}{0pt}{4pt}{1pt}

% ---------- theorem environments ----------
% NOTE ON ADAPTATION: this is an empirical study, not a theory report. It has no
% theorems. The epistemic-status apparatus is retained but repurposed: Hypothesis
% replaces Conjecture (a falsifiable prediction registered BEFORE the experiment),
% and Definition/Remark carry their usual load. Nothing here is proved, and
% nothing pretends to be.
\theoremstyle{definition}
\newtheorem{definition}{Definition}[section]
\newtheorem{hypothesis}{Hypothesis}[section]
\newtheorem{criterion}{Exit Criterion}[section]
\newtheorem{remark}{Remark}[section]

\crefname{hypothesis}{Hypothesis}{Hypotheses}
\crefname{definition}{Definition}{Definitions}
\crefname{criterion}{Exit Criterion}{Exit Criteria}
\crefname{section}{Section}{Sections}

% ---------- header/footer ----------
\setlength{\headheight}{12pt}
\pagestyle{fancy}\fancyhf{}
\fancyhead[L]{\scriptsize Cross-Geometry Generalization of DoMINO for External Aerodynamics}
\fancyhead[R]{\scriptsize Working Project Plan v1}
\fancyfoot[C]{\scriptsize\thepage}
\renewcommand{\headrulewidth}{0.3pt}
\renewcommand{\footrulewidth}{0pt}

% ---------- listings ----------
\lstset{
  basicstyle=\ttfamily\scriptsize,
  breaklines=true, frame=single, framesep=3pt, rulecolor=\color{black},
  language=Python, showstringspaces=false, columns=fullflexible,
  keywordstyle=\bfseries, commentstyle=\itshape,
  morekeywords={self,True,False,None}, aboveskip=4pt, belowskip=4pt,
}
\lstdefinestyle{sh}{language=bash,morekeywords={python,torchrun,git,pip}}
\lstdefinestyle{yml}{language=,morecomment=[l]{\#}}

% ---------- macros ----------
\DeclareMathOperator{\Cd}{C_d}
\newcommand{\R}{\mathbb{R}}
\newcommand{\calT}{\mathcal{T}}
\newcommand{\calS}{\mathcal{S}}
\newcommand{\norm}[1]{\lVert#1\rVert}
\newcommand{\abs}[1]{\lvert#1\rvert}
\newcommand{\Cp}{{C_p}}
\newcommand{\Cf}{{C_f}}

\title{\vspace{-1.5em}\large\bfseries
Cross-Geometry Generalization of DoMINO for External Aerodynamics\\
\normalsize Can a Neural Operator Trained on Simple Bluff Bodies Predict
Drag on a Real Car?\\[0.3em]
\small\mdseries Working Project Plan --- Phase-Wise Execution, Exit Criteria,
Metrics, Failure Modes, and Deliverables}
\author{}
\date{}

\begin{document}
\maketitle
\vspace{-3em}
\thispagestyle{fancy}

{\small
\noindent\textbf{Summary.}
We test whether the DoMINO neural operator, trained \emph{only} on simple
bluff-body aerodynamics data (AhmedML, WindsorML), can predict drag on a
realistic road car (DrivAerML) it has never seen --- and if not, how many
high-fidelity runs of the real car are required to recover. \Cref{sec:method} states
the learning problem in full: what operator DoMINO approximates, how its learnable
ball-query point-convolution kernels build a global geometry encoding, how a local
encoding is extracted per query point, and what loss is minimized. That mathematics is
not decoration --- it is what makes the central prediction falsifiable. Because the
convolution kernels are \emph{learned} from the training geometries, a body containing
features absent from the source distribution (wheels, mirrors, a notchback C-pillar)
has no learned representation, and error should localize \emph{there} and not on the
roof or hood. The DoMINO paper does test out-of-distribution samples, but along a
different axis --- drag force outside the training range, on morphs of the \emph{same}
body; this study tests transfer across geometric \emph{topology}
(\cref{rem:ood}). The headline deliverable is a single practically utilizable outcome:
\emph{drag coefficient $\Cd$ error on unseen geometry}, defended by two companion
metrics that exist to prevent a specific self-deception. A surrogate can obtain
approximately the right $\Cd$ by \emph{cancelling errors} --- over-predicting pressure
forward, under-predicting aft, and landing on a plausible integral. Such a model is
right for the wrong reasons and will fail the moment it is used to \emph{choose
between} designs. We therefore report, alongside $\Cd$: the Spearman rank correlation
across designs (does the model preserve design ordering?) and the area-weighted
surface $\Cp$ error map with signed bias (did it get $\Cd$ right honestly?).
Because a surface map can show \emph{where} the pressure is wrong but not \emph{why},
volume data is fetched for five cases only --- selected after the fact, once the error
ranking is known --- to evidence the mechanism rather than assert it
(\cref{sec:wake}). Predictions are registered as \emph{Hypotheses} \emph{before}
execution; the study is designed so that a negative result is the strongest outcome,
not a failure. Every
phase carries an \textbf{Exit Criterion} that must be cleared before proceeding ---
the dominant risks in this project are silent, and the exit criteria are what catch
them.
}
\vspace{4pt}

{\small
\noindent\textbf{Working with an AI assistant on this project?} Paste the operating
contract in \cref{sec:contract} (Appendix C) \emph{first}, before anything else. It
specifies what the assistant must verify, what it must refuse --- including requests
from you --- and what it should treat as suspicious. The dominant failure modes of this
study are silent, and an agreeable collaborator will help you commit them.
}
\vspace{4pt}

\tableofcontents
\vspace{6pt}
\hrule
\vspace{6pt}

% ============================================================
\section{Introduction}
% ============================================================

External aerodynamic simulation governs vehicle design: drag determines range,
fuel economy, and stability. High-fidelity CFD provides accurate predictions but
is computationally expensive, particularly for complex geometries, which places
a hard ceiling on how many designs can be evaluated in a development cycle. The
promise of machine-learned surrogates is to break that ceiling.

DoMINO (Ranade et al., NVIDIA) is a decomposable multi-scale iterative neural
operator that addresses the two failure modes of prior ML surrogates in this
domain --- limited accuracy and poor scalability, the latter typically forced by
aggressive mesh downsampling. DoMINO takes STL geometry as input, learns a
global multi-scale geometry encoding from the point cloud, extracts a local
geometry representation in a sub-region around each query point, and aggregates
that with a dynamically-constructed computational stencil to predict surface and
volume flow quantities. It is validated on DrivAerML and is demonstrated to be
mesh-independent.

The gap is precise, and it survives the paper's own generalization study. DrivAerML
comprises 500 parametrically morphed variants of \emph{one} vehicle --- a DrivAer
notchback. The paper does test out-of-distribution samples, but OOD there means
\emph{drag force outside the training range}, on morphs of that same body
(\cref{rem:ood}). Every geometry the model has ever seen has the same wheels, the same
mirrors, the same topology. Nothing in the literature answers the question every OEM
evaluating this technology actually has:

\begin{quote}
\emph{Will this model work on \textbf{my} car --- which is not a DrivAer
notchback?}
\end{quote}

This report specifies a study to answer that question, in the direction that
matters.

\subsection{The two directions of transfer, and why only one is worth compute}

``Unseen geometry'' admits two readings, with very different value.

\textbf{Downward transfer} (DrivAer $\to$ Ahmed/Windsor): train on the complex
car, test on the simple block. This will most likely succeed. It establishes
almost nothing --- generalizing from complex to simple is not a claim anyone
needs.

\textbf{Upward transfer} (Ahmed\,+\,Windsor $\to$ DrivAer): train on cheap simple
bodies, test on the expensive realistic car. This is the industrially meaningful
direction --- cheap data, expensive target --- and it is the direction in which
failure is \emph{informative}.

We take the upward direction as the experiment and retain the downward direction
as a control (\cref{sec:arms}).

\subsection{The single practical outcome}

\begin{definition}[Headline metric]
\label{def:headline}
The deliverable is $\Cd$ prediction error on unseen geometry, reported in
\emph{drag counts} ($1\ \text{count} = 0.001\ \Cd$), the unit in which
aerodynamicists actually speak.
\end{definition}

$\Cd$ is chosen over wall shear stress deliberately. Wall shear is small in
magnitude, noisy, dominated by near-wall behaviour the model never resolves, and
contributes only secondarily to drag on a bluff body, where pressure drag
dominates. It is the metric most likely to look poor for reasons unrelated to
the research question. It is retained as \emph{supporting evidence} --- it is
where separation-resolution failure becomes visible --- but not as the headline.

% ============================================================
\section{Background and Related Work}
% ============================================================

\subsection{The DoMINO architecture}

The model learns a global geometry encoding from the STL point cloud using a
multi-scale iterative approach that captures both short- and long-range
dependencies characteristic of elliptic PDEs, enriched with a signed distance
field (SDF) and positional encoding. Discrete points are sampled; a sub-region is
constructed around each; the local geometry encoding is extracted from the global
encoding via dynamic point convolution kernels. A computational stencil of $p+1$
points is then constructed dynamically by sampling neighbours within the
sub-region, and an aggregation network combines the local geometry encoding with
the stencil to predict solutions at the discrete points.

Critically for this study: the global geometry encoder is \emph{shared} between
surface and volume predictions, but the aggregation networks are \emph{separate},
because the surface variables (pressure, wall-shear vector) differ from the volume
variables (pressure, velocity, turbulence parameters). NVIDIA's own guidance is to
train surface and volume models separately. This project is surface-only.

\subsection{The dataset family}

The three automotive datasets were built by the same group with deliberately
consistent file structure and naming conventions, explicitly so that ML methods
can be tested \emph{across} them --- as CFD methods are validated on multiple
cases. This is the single largest piece of luck available to this project and the
reason it is tractable.

\begin{center}\small
\begin{tabular}{@{}p{2.2cm}p{1.3cm}p{3.6cm}p{2.6cm}p{4.2cm}@{}}
\toprule
\textbf{Dataset} & \textbf{Cases} & \textbf{Method} & \textbf{Body}
& \textbf{Role in this study} \\
\midrule
AhmedML   & 500 & Hybrid RANS-LES, $\sim$20M cells & Slanted-back block
  & Train (source); also \emph{sole} source in Arm A3a \\
WindsorML & 355 & WMLES                            & Squareback bluff body
  & Train (source) \\
DrivAerML & 500 & HRLES, $\sim$160M-cell volumes   & DrivAer notchback
  & \textbf{Target (held out)} \\
\midrule
DrivAerNet & 4000 & RANS, 8--16M cells (low-fidelity) & DrivAer
  & \emph{Not used} (see \cref{rem:scope}) \\
\bottomrule
\end{tabular}
\end{center}

\subsection{Positioning}

\begin{center}\small
\begin{tabular}{@{}p{3.6cm}p{3.0cm}p{3.4cm}p{5.2cm}@{}}
\toprule
\textbf{Work} & \textbf{Train} & \textbf{Test} & \textbf{Axis of generalization tested} \\
\midrule
DoMINO (Ranade et al.) & DrivAerML morphs & DrivAerML morphs, incl. OOD split
  & \emph{Drag-force range}, within one body. Also mesh-independence. \\[0.3em]
PhysicsNeMo-CFD benchmark & Various & Various
  & Aggregate accuracy / scalability across model architectures \\[0.3em]
\textbf{This work} & \textbf{Ahmed + Windsor} & \textbf{DrivAer}
  & \textbf{Geometric topology.} Unseen features: wheels, mirrors, C-pillar. \\
\bottomrule
\end{tabular}
\end{center}

\begin{remark}[What is and is not novel here --- state this honestly or be caught]
The DoMINO paper \emph{does} evaluate out-of-distribution samples, and a
PhysicsNeMo-CFD benchmarking framework already exists for assessing generalization of
AI models in automotive aerodynamics. ``Run DoMINO across datasets and tabulate'' is
therefore \textbf{not} a contribution, and claiming novelty for OOD testing as such
would be caught immediately by anyone who has read Section 4 of the paper.

The distinction --- and it is a real one --- is the \emph{axis}. Their OOD is defined by
drag force outside the training min--max range, evaluated on morphs of the same
notchback. Ours is a different body entirely. See \cref{rem:ood} for the full argument.
What remains unpublished, and what this study delivers, is (i) transfer across
\emph{geometric topology}, (ii) \emph{physics-resolved} failure localization --- where on
the body it breaks, and why --- and (iii) the \emph{sample-efficiency curve for recovery}.
The delta is the analysis and the transfer direction, not the training run.
\end{remark}

\begin{remark}[Scope discipline --- read before getting ambitious]
\label{rem:scope}
Deferred, deliberately: fidelity transfer (DrivAerNet RANS $\to$ DrivAerML HRLES),
uncertainty quantification, volume-field prediction, physics losses, and the
\emph{diversity-scaling} study (how many source body families does transfer require?
--- see \cref{rem:a3a}). Each roughly doubles runtime and halves the clarity of the
result. The reader of this artifact will spend approximately eight minutes on it and
must leave with \emph{one} sentence. One question, one number, three figures. Note also
that DrivAerNet is licensed CC-BY-NC (non-commercial), unlike the CC-BY/CC-BY-SA
licensing of the three datasets used here --- a further reason to exclude it from the
headline study.
\end{remark}

% ============================================================
\section{Method: What DoMINO Actually Learns}
\label{sec:method}

This section states the learning problem precisely. It exists because
\cref{hyp:degrade,hyp:localize} are claims about \emph{why the geometry encoder
fails to represent an unseen feature}. Those claims cannot be defended --- in a
review, in an interview, or to yourself --- without the mathematics they rest on.

\subsection{The learning problem}

\begin{definition}[The operator being approximated]
\label{def:operator}
Let $\Omega \subset \R^3$ be a computational domain and let $\calS \subset \Omega$
be the surface of a vehicle, presented as a triangulated STL. DoMINO approximates
the solution operator
\begin{equation}
\mathcal{G}: \calS \;\longmapsto\; \bigl(u,\, p,\, \nu_t \big|_\Omega;\;
p,\, \boldsymbol{\tau}_w \big|_{\calS}\bigr)
\label{eq:operator}
\end{equation}
mapping a \emph{geometry} to the flow fields it induces --- velocity, pressure and
turbulent viscosity in the volume; pressure and the wall-shear-stress vector on the
surface. This study concerns the surface branch only.
\end{definition}

\begin{remark}[Why ``neural operator'' and not ``neural network'']
The object learned is a map between \emph{function spaces}, not between fixed-dimensional
vectors. The practical consequence, and the reason it matters here, is that solutions
can be queried at \emph{arbitrary} points: the model is not tied to the mesh it was
trained on. The paper demonstrates this by evaluating surface predictions on a
uniform 10M-point cloud sampled from the STL instead of the simulation mesh, and
recovering the same drag regression ($R^2 = 0.96$ in both cases). This
\emph{mesh-independence} is the property that makes the architecture industrially
interesting --- meshing is itself expensive --- and it is why transfer to a new body is
even a coherent question to ask.
\end{remark}

\begin{remark}[What this is, stripped of branding]
$\mathcal{G}$ is approximated by supervised regression against CFD ground truth,
trained with MSE and Adam. There is no reasoning, no physical model, and no
understanding of fluid dynamics inside DoMINO. It is a \textbf{learned interpolator
over a geometry-to-field mapping}. This is not a criticism; it is the fact that makes
this project's question well-posed. A system that \emph{understood} aerodynamics would
generalize to a wheel because it knows what a wheel does to flow. An interpolator
generalizes to a wheel only if something wheel-like was inside its training
distribution. \Cref{hyp:degrade} follows directly from this observation.
\end{remark}

\subsection{Stage 1 --- Global geometry representation}

Two bounding boxes are constructed: a tight-fitted box around the surface, and a
larger box representing the computational domain. In the reference configuration the
domain extends 2 car lengths downstream, 1 car length upstream, and 1 car
width/height in the transverse directions.

The STL point cloud is projected onto a structured grid of resolution
$m\times m\times m\times f$ using \emph{learnable point convolution kernels}. This is
the fundamental primitive of the architecture:

\begin{equation}
\boxed{\;
y_i = \sum_{j=0}^{n_y} f\bigl(\vec{x}_i,\ \vec{x}_j,\ d_{ij}\bigr)
\;}
\label{eq:pointconv}
\end{equation}

where $\vec{x}_i, \vec{x}_j$ are points from two point clouds, $d_{ij}$ is the
distance between them, and $f$ is a fully connected network. The sum runs over $n_y$
neighbours found by a \textbf{ball query} --- implemented as custom GPU kernels in
NVIDIA Warp.

\begin{definition}[Radius of influence]
\label{def:radius}
\Cref{eq:pointconv} is governed by two parameters: the number of neighbours $n_y$,
and the \textbf{radius of influence} $r$, which sets the physical size of the kernel.
A \emph{large} $r$ propagates geometry information further into the domain and
captures long-range interaction; a \emph{small} $r$ captures fine geometric detail.
DoMINO is \emph{multi-scale} precisely in this sense: it employs a \emph{range} of
radii simultaneously, so that both fine features and long-range dependence are
represented.
\end{definition}

Geometry information reaches the computational-domain grid $G_c$ by two routes:
\begin{enumerate}[label=(\roman*)]
  \item a separate set of multi-scale point convolution kernels projecting directly
  onto $G_c$; and
  \item \textbf{iterative} CNN blocks (convolution, pooling, unpooling) propagating
  features from the surface grid $G_s$ to $G_c$, applied for a fixed number of
  iterations. (The paper notes a fixed-point stopping criterion such as
  $L_2(G_s, G_c) < 10^{-3}$ is possible but was not explored.) This is the
  ``\textbf{Iterative}'' in DoMINO.
\end{enumerate}

Finally, the \textbf{signed distance field} (SDF) and its gradient components are
computed on $G_c$ and \emph{appended} to the learned features, supplying explicit
topological information about the geometry. The result is a dense, high-dimensional
\textbf{global encoding} of the body.

\subsection{Stage 2 --- Local geometry representation}

\begin{remark}[The central architectural insight]
The global encoding is dense and high-dimensional, but the solution at any point is
dominated by information in its \emph{locality}. Using the global encoding directly
is therefore both wasteful and inaccurate. DoMINO instead \emph{extracts} a local
encoding per query point. This is the ``\textbf{Decomposable}'' in the name, and it
is what allows the model to scale to $\sim$150M-element meshes without downsampling
--- the failure mode that the paper identifies in prior ML surrogates.
\end{remark}

For each sampled point, $p$ neighbours are drawn to form a computational stencil of
$p+1$ points (by analogy with finite-volume/finite-element stencils). A sub-region
is drawn around the stencil, and a point convolution kernel (\cref{eq:pointconv})
extracts local features of resolution $l\times l\times l\times n_f$ from the global
encoding, which are then transformed by fully connected layers.

\subsection{Stage 3 --- Aggregation network}

Each point in the stencil carries input features:
\begin{itemize}
  \item physical coordinates in the computational domain;
  \item the SDF at those coordinates;
  \item the normal vector from the centre of mass of the domain;
  \item the \emph{surface} normal --- only if the point lies on the surface.
\end{itemize}

These pass through a \textbf{basis function neural network} producing a latent vector
per stencil point. Each latent vector is concatenated with the local geometry
encoding, passed through further fully connected layers to give a solution vector at
each stencil point, and the $p+1$ predictions are combined by \textbf{inverse-distance
weighting} to yield the solution at the sampled point. A separate aggregation network
instance is used per solution variable; the local geometry representation is shared
across them.

\subsection{The loss function}

\begin{definition}[Training objective]
\label{def:loss}
For surface variables $y \in \{p, \tau_x, \tau_y, \tau_z\}$, the loss is MSE on
sampled mesh points, computed \emph{per variable} and summed, plus an
\textbf{area-weighted} MSE term on the surface predictions:
\begin{equation}
\mathcal{L}
= \underbrace{\sum_{v} \frac{1}{|\mathcal{P}|}\sum_{i \in \mathcal{P}}
   \bigl(y^{pred}_{v,i} - y^{true}_{v,i}\bigr)^2}_{\text{pointwise MSE, per variable}}
\;+\;
  \underbrace{\sum_{v} \frac{1}{|\mathcal{P}|}\sum_{i \in \mathcal{P}}
   \bigl(a_i\, y^{pred}_{v,i} - a_i\, y^{true}_{v,i}\bigr)^2}
   _{\text{area-weighted MSE}}
\label{eq:loss}
\end{equation}
where $\mathcal{P}$ is the subset of points sampled this epoch and $a_i$ the facet
area. Optimizer: Adam, reduce-on-plateau from $10^{-3}$ to $10^{-6}$, 500 epochs,
AMP float-16.
\end{definition}

\begin{remark}[The variable-order assumption originates here --- VERIFIED IN SOURCE]
\Cref{eq:loss} is summed \emph{per variable}, and the surface loss implementations index
into the target tensor \emph{positionally}. This is not inferred from documentation; it
was read from the code. In
\texttt{examples/cfd/external\_aerodynamics/domino/src/loss.py}, \texttt{drag\_loss\_fn}
(lines 423--427 at commit \texttt{59aaf59}) contains:
\begin{lstlisting}
pres_true = output_true[:, :, 0] * normals[:, :, 0]   # channel 0 IS pressure
wx_true   = output_true[:, :, 1]                       # channel 1 IS x-wall-shear
\end{lstlisting}
Channel 0 \emph{must} be pressure; channel 1 \emph{must} be $x$-wall-shear. Nothing
validates this --- the tensor carries no names, so nothing \emph{can}. A permuted target
still produces a well-formed, smoothly decreasing loss; it simply regresses the wrong
channels against each other. This is the mechanical origin of the trap in
\cref{rem:ordering}, and it is now a fact rather than a warning.

\textbf{Two further consequences, not previously noted.} (i) Only channel 1 enters the
drag loss --- $\tau_y$ and $\tau_z$ are absent, consistent with drag being streamwise, but
it means channels 2 and 3 could be swapped without the \emph{drag} loss noticing (though
\texttt{loss\_fn\_surface} would). (ii) \texttt{normals[:, :, 0]} encodes an assumption
that \textbf{$x$ is the streamwise axis} --- see \cref{rem:streamwise}.
\end{remark}

\begin{remark}[The streamwise axis is assumed to be $x$ --- verify per dataset]
\label{rem:streamwise}
\texttt{drag\_loss\_fn} multiplies pressure by \texttt{normals[:, :, 0]}, the
$x$-component of the surface normal. The drag loss therefore assumes the freestream is
aligned with $+x$ in \emph{every} dataset. Ahmed, Windsor, and DrivAer were built by the
same group and very likely agree --- but ``very likely'' is how silent failures are
introduced. If any body is oriented differently, the integral loss will optimize a
component of force that is not drag, and, as with the ordering trap, it will do so
without complaint.

Add a row to the compatibility table (\cref{sec:phase0}): \textbf{which axis is
streamwise?} Confirm from the STL bounding box and the published force conventions, not
from assumption.
\end{remark}

\begin{remark}[What the tree actually contains --- read at commit \texttt{59aaf59}]
\label{rem:treereality}
\Cref{rem:structstatus} warned that the layout below was taken from documentation rather
than inspection. It has now been inspected, and \textbf{the documentation is stale
relative to its own repository}. Recorded here because the corrections are load-bearing.

\begin{center}\small
\begin{tabular}{@{}p{4.2cm}p{10.8cm}@{}}
\toprule
\textbf{Plan / README claimed} & \textbf{Reality at \texttt{59aaf59}} \\
\midrule
\texttt{openfoam\_datapipe.py}
  & \textbf{Does not exist.} The datapipe is
    \texttt{physicsnemo/datapipes/cae/domino\_datapipe.py} --- \emph{library} code, not
    example code. Preprocessing has moved to PhysicsNeMo-Curator. \\[0.3em]
\texttt{retraining.py}
  & \textbf{Does not exist.} Folded into \texttt{train.py} via \texttt{resume\_dir}.
    See \cref{rem:epochtrap} --- using it naively silently breaks the recovery curve. \\[0.3em]
\texttt{process\_data.py}
  & \textbf{Does not exist.} \\[0.3em]
(not mentioned in the plan)
  & \texttt{src/loss.py} \emph{does} exist, and is where the ordering trap lives
    (\cref{rem:ordering}). \\[0.3em]
(not mentioned in the plan)
  & \texttt{unified\_external\_aero\_recipe/} exists, with its own
    \texttt{domino\_surface.yaml}. There may now be two supported paths; this study uses
    the \texttt{examples/.../domino} example. \\[0.3em]
\texttt{train.py}, \texttt{test.py},
\texttt{compute\_statistics.py},
\texttt{cache\_data.py},
\texttt{conf/config.yaml},
\texttt{conf/cached.yaml}
  & All present, under \texttt{examples/cfd/external\_aerodynamics/domino/src/}. \\
\bottomrule
\end{tabular}
\end{center}

The example's own README instructs the user to modify \texttt{openfoam\_datapipe.py} and
to run \texttt{process\_data.py}. Neither exists. The README is internally inconsistent
with the code it ships alongside. \textbf{Trust the tree, not the documentation} ---
including this document, which is why it is being corrected rather than defended.
\end{remark}

\subsection{Published results, and what they do and do not establish}
\label{sec:published}

\begin{center}\small
\begin{tabular}{@{}lcc@{}}
\toprule
\textbf{Surface field} & \textbf{Rel. $L_2$} & \textbf{Area-weighted rel. $L_2$} \\
\midrule
Pressure       & 0.1505 & 0.1181 \\
X-wall-shear   & 0.2124 & 0.1279 \\
Y-wall-shear   & 0.3020 & 0.2769 \\
Z-wall-shear   & 0.3359 & 0.2290 \\
\bottomrule
\end{tabular}
\end{center}

with relative error defined as
\begin{equation}
\epsilon = \frac{\sqrt{\sum \left(y_T - y_P\right)^2}}{\sqrt{\sum y_T^2}}
\label{eq:relerr}
\end{equation}
Drag regression: $R^2 = 0.96$, on both the simulation mesh and a uniform 10M-point
cloud. Area-weighted errors are uniformly \emph{lower} than raw ones, indicating
error concentrates on small-area facets; $y$- and $z$-wall-shear are worst, owing to
their small magnitudes.

\begin{remark}[The paper already tests OOD --- but a different OOD. Read this carefully.]
\label{rem:ood}
DoMINO \textbf{is} evaluated out-of-distribution. The test split reserves 10\% of
samples, of which roughly 20\% are OOD, and OOD is defined by \emph{drag force falling
outside the training min--max range}; the showcased cases (IDs 419, 439) are the
lowest- and highest-drag designs in the dataset, differing by $\sim$300\,N.

This is \textbf{out-of-distribution in the output range, within a single body}. Every
sample --- train and test, ID and OOD --- is a morph of the same DrivAer notchback, with
the same wheels, the same mirrors, the same topology.

What this study tests is \textbf{out-of-distribution in geometric topology}: a body
whose features the encoder has never seen in any form. These are different axes of
generalization and must be distinguished explicitly in the writeup, or a reader who
knows the paper will conclude the work has been redone.

The paper's result in fact \emph{sharpens} this study rather than pre-empting it.
DoMINO demonstrably handles OOD drag \emph{values} well. If it fails on OOD
\emph{geometry}, the contrast isolates precisely which axis of generalization breaks
--- a cleaner finding than either result alone.
\end{remark}

\subsection{Why unseen geometry should be expected to break it}
\label{sec:whybreak}

The mechanism can now be stated exactly, and it is the load-bearing argument of the
whole study.

\begin{enumerate}[label=(\arabic*)]
  \item The point convolution kernels of \cref{eq:pointconv} are \textbf{learned}.
  The function $f$ is fit to the training geometries. Its parameters encode which
  local point-cloud configurations matter and how they map to flow features.
  \item The local encoding is likewise learned: the paper states explicitly that
  local geometry representations \emph{extract the essential information required to
  predict solution fields in different regions} --- e.g.\ ``in predicting the flow
  fields in the wake, only the local geometry features required for learning these are
  extracted.'' Which features are ``essential'' is \emph{determined by the training
  distribution}.
  \item Ahmed and Windsor contain no wheel, no wheelhouse, no mirror, and no
  notchback C-pillar. No configuration of points resembling a rotating-wheel
  wheelarch, or a mirror stalk shedding a vortex, ever passes through $f$ during
  training.
  \item Consequently the kernels have no basis on which to encode them, and the
  aggregation network has never been asked to map such a local encoding to a
  surface pressure. The SDF will faithfully report that geometry \emph{is} there ---
  but a faithful SDF of an unseen feature is not the same as a learned representation
  of what that feature does to the flow.
\end{enumerate}

This is the content of \cref{hyp:degrade} and \cref{hyp:localize}. Note that it makes
a \emph{spatially localized} prediction --- error should concentrate on the specific
features absent from the source bodies, and \emph{not} on the roof or hood, which
Ahmed and Windsor do resemble. That is a falsifiable claim, and Phase 5 is built to
adjudicate it.

\begin{remark}[The optimistic counter-argument, stated fairly]
It could go the other way, and the reason is worth understanding. DoMINO is
\emph{local} by construction: predictions depend on a sub-region and a $p+1$-point
stencil, not on the global body. If the local flow physics around a mirror is
approximately the local flow physics around \emph{some} bluff-body feature the model
has seen --- a sharp edge, a separation line, a slanted surface --- then locality could
carry the model further than the global-topology argument suggests. Ahmed's slanted
rear does produce geometry-induced separation; Windsor's squareback does produce a
base wake. The question is empirical, which is exactly why it is worth the compute.
\end{remark}

\subsection{Notation summary}

\begin{center}\small
\begin{tabular}{@{}llp{9.6cm}@{}}
\toprule
\textbf{Symbol} & \textbf{Object} & \textbf{Meaning} \\
\midrule
$\mathcal{G}$ & Operator & Geometry $\mapsto$ flow fields (\cref{eq:operator}) \\
$\calS,\ \Omega$ & Sets & Vehicle surface; computational domain \\
$G_s,\ G_c$ & Grids & Surface bounding-box grid; computational-domain grid \\
$m \times m \times m \times f$ & Resolution & Structured global encoding \\
$l \times l \times l \times n_f$ & Resolution & Extracted local encoding \\
$r$ & Scalar & Ball-query radius of influence (\cref{def:radius}) \\
$n_y$ & Integer & Neighbours in a point convolution kernel \\
$p$ & Integer & Stencil neighbours; stencil has $p+1$ points \\
$d_{ij}$ & Scalar & Distance metric between points $i$ and $j$ \\
$a_i$ & Scalar & Facet area, used in area-weighted loss and integration \\
$\boldsymbol{\tau}_w$ & Vector & Wall-shear-stress \\
\bottomrule
\end{tabular}
\end{center}

% ============================================================
\section{The Core Question and Registered Hypotheses}
\label{sec:hypotheses}
% ============================================================

\subsection{Statement}

\begin{quote}
\emph{If a DoMINO surrogate is trained only on simple bluff bodies (Ahmed and
Windsor), can it predict the drag of a realistic road car (DrivAer) it has never
seen --- and if not, how many high-fidelity runs of the real car does it take to
fix?}
\end{quote}

\subsection{Registered hypotheses}

These are stated \emph{before} execution and are falsifiable. They are labelled
Hypotheses, not results. \textbf{Nothing in this document is proved.}

\begin{hypothesis}[Upward transfer degrades]
\label{hyp:degrade}
A DoMINO model trained on AhmedML + WindsorML will show materially worse $\Cd$
error on DrivAerML than a model trained on DrivAerML itself. Mechanism
(\cref{sec:whybreak}): the point convolution kernels of \cref{eq:pointconv} are
\emph{learned} from the training geometries. No configuration of points resembling a
wheelarch, a mirror stalk, or a notchback C-pillar ever passes through $f$ during
training on Ahmed and Windsor, so the kernels have no basis on which to encode them.
The SDF will report that geometry is present; it will not supply a learned
representation of what that geometry does to the flow.
\end{hypothesis}

\begin{hypothesis}[Error localizes to absent features]
\label{hyp:localize}
Under cold upward transfer, surface $\Cp$ error will concentrate on the wheels,
wheelhouses, mirrors, and C-pillar/rear-window region --- precisely the features
absent from the source bodies --- while simple zones (roof, hood) remain
comparatively accurate.
\end{hypothesis}

\begin{hypothesis}[Error cancellation precedes accuracy]
\label{hyp:cancel}
At small fine-tuning budgets, the model will achieve plausible $\Cd$ before it
achieves an accurate $\Cp$ field: area-weighted $\Cp$ L2 error will remain large
while the \emph{signed} bias approaches zero, indicating large local errors that
cancel in the integral. Design-ranking correlation will lag absolute $\Cd$
accuracy.
\end{hypothesis}

\begin{hypothesis}[Downward transfer succeeds]
\label{hyp:downward}
A model trained on DrivAerML will transfer \emph{successfully} to Ahmed and
Windsor. This is the control: its success is what forecloses the objection that a
failure in \cref{hyp:degrade} merely reflects a broken pipeline.
\end{hypothesis}

\begin{hypothesis}[Source diversity buys transfer]
\label{hyp:diversity}
A model trained on Ahmed + Windsor will show materially lower cold-transfer $\Cd$
error on DrivAer than a model trained on Ahmed alone. This is an
\emph{interpretive control}, not a scaling law: with two points it cannot say how many
body families are needed, only whether adding one moves the number at all. If it does
not move, the barrier is topological rather than one of source-set size --- a stronger
reading of \cref{hyp:degrade}.
\end{hypothesis}

\begin{remark}[Why a negative result is the strong outcome]
The study is constructed so that it cannot fail into worthlessness. If upward
transfer succeeds, cheap bluff-body data buys real-car predictions --- a cost
argument with money attached. If it fails, the limits of geometric generalization
in neural operators are demonstrated, and the recovery curve states exactly what
repair costs. \Cref{hyp:degrade} is expected to hold. That expectation is the
point, not a hedge.
\end{remark}

\subsection{The three reported quantities}

\begin{definition}[Accuracy]
$\Cd$ mean absolute error in drag counts on held-out DrivAer designs.
\end{definition}

\begin{definition}[Usability]
\label{def:rank}
Spearman rank correlation $\rho$ between predicted and true $\Cd$ across designs.
A surrogate that cannot order two designs by drag is useless for optimization
\emph{regardless} of its absolute error. This metric, not $\Cd$ MAE, determines
whether the model is deployable in a design loop.
\end{definition}

\begin{definition}[Honesty]
\label{def:honesty}
Area-weighted surface $\Cp$ error and its \emph{signed} bias. Given per-face
areas $a_i$ with $w_i = a_i / \sum_j a_j$ and error $e_i = \Cp^{pred}_i -
\Cp^{true}_i$:
\begin{equation}
\mathcal{E}_{L2} = \sqrt{\textstyle\sum_i w_i e_i^2},
\qquad
\mathcal{E}_{bias} = \textstyle\sum_i w_i e_i
\label{eq:honesty}
\end{equation}
\end{definition}

\begin{remark}[$\mathcal{E}_{bias}$ is the error-cancellation detector]
\label{rem:cancel}
If $\mathcal{E}_{L2}$ is large while $\mathcal{E}_{bias} \approx 0$, the model is
making large local errors that cancel under integration --- obtaining the right
$\Cd$ for the wrong reasons. This is exactly the pathology that makes a surrogate
dangerous in an optimization loop, and it is invisible to $\Cd$ alone. Log both,
every run, without exception.
\end{remark}

% ============================================================
\section{Nondimensionalization}
\label{sec:nondim}
% ============================================================

The three bodies differ in frontal area, length scale, and Reynolds number. A
consistent nondimensionalization is not a formatting preference; it is a
correctness precondition.

\begin{equation}
\Cp = \frac{p - p_\infty}{\tfrac{1}{2}\rho U_\infty^2},
\qquad
\Cf = \frac{\boldsymbol{\tau}_w}{\tfrac{1}{2}\rho U_\infty^2},
\qquad
\Cd = \frac{F_x}{\tfrac{1}{2}\rho U_\infty^2 A_{\text{frontal}}}
\label{eq:nondim}
\end{equation}

\begin{definition}[Canonical output space]
\label{def:canonical}
Every body, every dataset, one target vector, one fixed order:
\begin{equation}
\boxed{\;
\mathbf{y} = \bigl[\, \Cp,\; \Cf_{,x},\; \Cf_{,y},\; \Cf_{,z} \,\bigr]
\;}
\label{eq:canonical}
\end{equation}
nondimensionalized by that body's own $\rho$ and $U_\infty$. The ordering is
\emph{not} arbitrary --- see \cref{rem:ordering}.
\end{definition}

\begin{remark}[The variable-ordering trap --- the single most dangerous item in this document]
\label{rem:ordering}
The DoMINO surface loss functions (\texttt{integral\_loss\_fn},
\texttt{loss\_fn\_surface}, \texttt{loss\_fn\_area}) \textbf{assume the variables
are ordered pressure-first, then the wall-shear-stress vector.} The NVIDIA
documentation states this explicitly and instructs users to modify the losses if
their ordering differs. If any of the three datasets stores them differently and
this goes undetected, the loss will be computed against the wrong targets and
\emph{training will appear to work perfectly}: a smooth, converging loss curve
producing a worthless model. This failure is silent, and it will not surface until
the results are inexplicable. Verify ordering per dataset in \cref{sec:phase0};
enforce \cref{eq:canonical} in the datapipe; assert it.
\end{remark}

\begin{remark}[Reference area must be declared, not assumed]
Decide and commit, in writing, how $A_{\text{frontal}}$ is computed for each body
(projected area of the STL onto the $y$--$z$ plane is the usual convention) and
whether it matches the reference area used by the dataset authors when they
published their own force coefficients. If these disagree, every $\Cd$ in the study
is off by a constant factor and the error will be found by the first competent
reader.
\end{remark}

% ============================================================
\section{Phase 0 --- Reconnaissance and Compatibility}
\label{sec:phase0}
% ============================================================

\textbf{Goal:} eliminate every unknown capable of silently corrupting a later
phase. No GPU work occurs here. Budget: 1--2 sessions.

\subsection{Environment and repository}

Create the project structure specified in \cref{sec:structure} \emph{first} --- the
two-repository split (\cref{rem:twolevel}) and the scaling-factor location
(\cref{rem:scalingpath}) are much cheaper to establish now than to retrofit.

\begin{lstlisting}[style=sh]
mkdir -p ~/domino-xgeom/{external,patches,conf,src,tests,results,figures}
cd ~/domino-xgeom && git init

cd external
git clone https://github.com/NVIDIA/physicsnemo.git          # model + training
git clone https://github.com/NVIDIA/physicsnemo-curator.git  # preprocessing
git clone https://github.com/NVIDIA/physicsnemo-cfd.git      # force-util oracle
cd physicsnemo && git rev-parse HEAD    # -> record in environment.md
\end{lstlisting}

Record PhysicsNeMo commit hash, CUDA version, driver, GPU count and type in
\texttt{environment.md} on day one. Reproducibility is part of the deliverable, not an
afterthought --- and ``which release?'' is the first question anyone asks when numbers
fail to reproduce.

\textbf{Then verify \cref{sec:structure} against the actual tree.} The internal layout
of the DoMINO example directory in that appendix is taken from documentation, not from
inspection (\cref{rem:structstatus}). Correct it now, while it is cheap.

\subsection{Storage}

DrivAerML volume solutions run to $\sim$150M elements per case; the full dataset is
tens of TB. \textbf{The volumes are not needed} (\cref{rem:novtu}). Download STL and
surface VTP; skip VTU. Reserve \texttt{drivaer\_holdout/} now and do not touch it until
Phase 5.

\subsection{Reconnaissance script}

\begin{lstlisting}
"""recon.py -- verify dataset compatibility BEFORE training anything."""
import pyvista as pv, numpy as np, json, sys
from pathlib import Path

def inspect(case_dir: Path) -> dict:
    stl  = next(case_dir.glob("*.stl")); mesh = pv.read(stl)
    vtp  = next(case_dir.glob("*.vtp")); surf = pv.read(vtp)
    b = mesh.bounds                      # (xmin,xmax,ymin,ymax,zmin,zmax)
    return {
        "case": case_dir.name,
        "stl_n_cells": mesh.n_cells,
        "bounds": list(b),
        "L_x": b[1]-b[0], "W_y": b[3]-b[2], "H_z": b[5]-b[4],
        # THE CRITICAL ONES: field names, and the ORDER they appear in.
        "point_arrays": list(surf.point_data.keys()),
        "cell_arrays":  list(surf.cell_data.keys()),
        # value ranges -- catches unit mismatch (Pa vs already-normalised Cp)
        "field_ranges": {k: [float(np.min(surf[k])), float(np.max(surf[k]))]
                         for k in surf.point_data.keys()},
    }

if __name__ == "__main__":
    root = Path(sys.argv[1])
    cases = sorted(d for d in root.iterdir() if d.is_dir())[:5]
    out = [inspect(c) for c in cases]
    Path(f"recon_{root.name}.json").write_text(json.dumps(out, indent=2))
    print(json.dumps(out, indent=2))
\end{lstlisting}

\subsection{The compatibility table}

Completed from \emph{inspected data} (\texttt{src/recon.py}) and the three dataset
papers. Not from documentation --- NVIDIA's README has already been shown to
contradict its own repository (\cref{rem:treereality}).

\begin{center}\footnotesize
\begin{tabular}{@{}p{3.5cm}p{3.4cm}p{3.4cm}p{3.4cm}@{}}
\toprule
\textbf{Question} & \textbf{Ahmed} & \textbf{Windsor} & \textbf{DrivAer} \\
\midrule
Surface file & \texttt{boundary\_*.vtp} & \textbf{\texttt{boundary\_*.vtu}} & \texttt{boundary\_*.vtp} \\
Fields on points or cells? & \textbf{cell} & \textbf{point} & \textbf{cell} \\
$\Cp$ field name & \texttt{static(p)\_coeffMean} & \texttt{cpavg} & \texttt{CpMeanTrim} \\
Dimensional $p$ field & \texttt{pMean} & \textbf{none published} & \texttt{pMeanTrim} \\
Wall-shear storage & one 3-vector, \textbf{Pa} & \textbf{three scalars}, already $\Cf$ & one 3-vector, \textbf{Pa} \\
Wall-shear name(s) & \texttt{wallShearStressMean} & \texttt{cfxavg},\ \texttt{cfzavg},\ \texttt{cfyavg} \emph{(file order: z before y)} & \texttt{wallShearStress\-MeanTrim} \\
$U_\infty$ & \textbf{1 m/s} & 30 m/s & 38.889 m/s \\
$\nu$ & $3.75\times10^{-7}$ & --- & $1.507\times10^{-5}$ \\
Constant $A_{\text{ref}}$ & 0.112 m$^2$ \emph{(paper only)} & 0.112 m$^2$ & 2.17 m$^2$ \\
\textbf{Constant-ref force file} & \texttt{force\_mom\_*.csv} & \texttt{force\_mom\_*.csv} & \textbf{\texttt{force\_mom\_constref\_*}} \\
Variable-ref force file & \texttt{force\_mom\_varref\_*} & \texttt{force\_mom\_varref\_*} & \textbf{\texttt{force\_mom\_*.csv}} \\
$\Cd$ column header & \texttt{cd} \emph{(lowercase)} & \texttt{cd} \emph{(lowercase)} & \texttt{Cd} \emph{(uppercase)} \\
\textbf{Streamwise axis} & \textbf{x} & \textbf{x} & \textbf{x} \\
\textbf{Vertical axis} & \textbf{z} & \textbf{y} & \textbf{z} \\
\textbf{Yaw} & $0^\circ$ & $\mathbf{-2.5^\circ}$ & $0^\circ$ \\
\bottomrule
\end{tabular}
\end{center}

\begin{remark}[Every column of that table contains a trap]
\label{rem:traps}
The three datasets were published by the same group, in a deliberately consistent
format, explicitly so that ML methods could be tested across them. They are
nonetheless incompatible in seven ways, \emph{all} of which fail silently. The table
above is the output of \texttt{src/recon.py} and the three dataset papers; it is
encoded, with provenance, in \texttt{src/datasets.py}, which is the single source of
truth. \textbf{Nothing else in this project may hardcode a field name or a reference
area.} The individual traps are stated in
\cref{rem:refarea,rem:windsorfmt,rem:velocity,rem:windsoraxes}.
\end{remark}

\begin{remark}[The reference-area trap: \texttt{force\_mom\_1.csv} inverts meaning]
\label{rem:refarea}
Each dataset publishes \emph{two} force files: one normalized by a constant reference
area, one by each case's own frontal area. The filenames are the trap.

\begin{center}\small
\begin{tabular}{@{}lll@{}}
\toprule
& \textbf{\texttt{force\_mom\_<i>.csv}} & \textbf{Suffixed file} \\
\midrule
Ahmed   & \textbf{constant} & \texttt{\_varref} $=$ variable \\
Windsor & \textbf{constant} & \texttt{\_varref} $=$ variable \\
DrivAer & \textbf{variable} & \texttt{\_constref} $=$ constant \\
\bottomrule
\end{tabular}
\end{center}

The unsuffixed file means \emph{opposite things} in the source datasets and the target.
Loading \texttt{force\_mom\_1.csv} uniformly --- the obvious implementation --- yields
constant-reference ground truth from Ahmed and Windsor and variable-reference ground
truth from DrivAer, across \emph{precisely} the boundary this study transfers over.

\textbf{The magnitude is not marginal.} For Ahmed run\_1 the two conventions differ by
\textbf{47 drag counts} (0.2385 vs 0.2854); for DrivAer run\_1, by 7.4. The headline
metric of this study is measured in drag counts, and the effect being measured is
smaller than the error this would introduce.

Confirmed from three independent sources: the AhmedML paper (Table 3), the WindsorML
paper (which states the convention explicitly), and DrivAerML's own
\texttt{geo\_ref\_<i>.csv}, which publishes \texttt{aRef} $=2.223$ (this morph) alongside
\texttt{aRefRef} $=2.17$ (constant) --- and $0.3035 \times (2.223/2.17) = 0.3109$,
reconciling the two files arithmetically.

\textbf{This study uses the constant-reference file throughout}, which means
\texttt{force\_mom\_*.csv} for Ahmed and Windsor but
\texttt{force\_mom\_constref\_*.csv} for DrivAer. The reason is \cref{def:rank}: the
usability metric is a design \emph{ranking}, and under a per-case area a morph that
changes frontal area moves its own $\Cd$ for reasons unrelated to the flow.
\end{remark}

\begin{remark}[Windsor's surface is a \texttt{.vtu} with point data]
\label{rem:windsorfmt}
Ahmed and DrivAer ship \texttt{boundary\_<i>.vtp} --- polydata, fields on \emph{cells}.
Windsor ships \texttt{boundary\_<i>.vtu} --- an unstructured grid, fields on
\emph{points}. It is a surface in a volume container.

Two consequences. A loader globbing for \texttt{*.vtp} finds nothing in Windsor. And
area-weighted integration (\cref{eq:forces}) requires \emph{cell} data, so Windsor must
be converted before any force is computed --- \texttt{point\_data\_to\_cell\_data()}, not
optional.

Note also that PhysicsNeMo-Curator ships worked examples for \texttt{ahmedml} and
\texttt{drivaerml} but \textbf{not} for Windsor. Windsor is therefore the custom dataset,
and the bulk of this project's engineering delta lives there.
\end{remark}

\begin{remark}[$U_\infty$ spans 39$\times$, so wall shear spans $\sim$1500$\times$]
\label{rem:velocity}
Ahmed runs at $U_\infty = \mathbf{1}$ m/s. Not 30. The AhmedML paper sets its Reynolds
number ($7.68\times10^5$ on body height) through the \emph{viscosity},
$\nu = 3.75\times10^{-7}$, not the velocity. Windsor runs at 30 m/s; DrivAer at 38.889.

Wall shear scales as $U_\infty^2$, so the \emph{dimensional} shear fields differ by a
factor of roughly 1500 between Ahmed and DrivAer. This is directly visible in the recon
output: Ahmed's \texttt{wallShearStressMean} spans $\pm0.036$; DrivAer's
\texttt{wallShearStressMeanTrim} spans $\pm40$.

Training on both without non-dimensionalizing would let DrivAer dominate the loss
entirely --- Ahmed's contribution would be numerical noise. This is not a subtle
weighting issue; it is three orders of magnitude.

Windsor is the exception: it publishes $\Cf$ directly, already non-dimensional. Ahmed and
DrivAer must be divided by $q = \tfrac{1}{2}\rho U_\infty^2$. \textbf{None of the three
datasets publishes $\rho$} --- they are incompressible simulations specified by kinematic
viscosity. Phase 2 must verify the assumed $\rho$ by recomputing $\Cp$ from $p$ and $q$
and checking against each dataset's own published $\Cp$ field. If they disagree, $\rho$
is wrong.
\end{remark}

\begin{remark}[Windsor is y-up and yawed; drag survives, lift does not]
\label{rem:windsoraxes}
The WindsorML paper: \emph{``for all the simulations in this work, y is upwards, which
differs to the original Windsor geometry where z is upwards.''} Recon confirms it ---
Windsor's STL spans $y \in [0, 0.475]$ (height) and $z \in [-0.194, +0.194]$ (width).
Ahmed and DrivAer are z-up.

\textbf{Drag is safe.} $x$ is streamwise in all three (verified from the STL bounding
boxes), and drag depends only on the $x$-component. \Cref{rem:streamwise} is therefore
resolved favourably: DoMINO's \texttt{normals[:, :, 0]} assumption holds.

\textbf{Lift, side force, and the $y$/$z$ shear channels are not.} The canonical
$\Cf_{,y}$ and $\Cf_{,z}$ mean \emph{different physical directions} in Windsor than in
the other two. The model is trained on all four channels, so a permuted $y$/$z$ in one
third of the training data is exactly the silent corruption this project is built to
prevent. \textbf{Phase 3 must decide explicitly whether to swap Windsor's axes into the
Ahmed/DrivAer convention.}

Compounding this: Windsor's wall shear is stored as three separate scalars in file order
\texttt{cfxavg}, \texttt{cfzavg}, \texttt{cfyavg} --- \emph{z before y}. Stacking them in
array order produces a permuted vector \emph{on top of} the axis difference.

\textbf{And Windsor is yawed $-2.5^\circ$}, generating a side force by design; Ahmed and
DrivAer are at zero yaw. Windsor's flow is asymmetric where the others are symmetric.
This is a real physical difference between the source bodies and the target, and belongs
in the limitations.
\end{remark}

\begin{criterion}[Phase 0 --- CLEARED]
\label{crit:p0}
\textbf{Cleared 2026-07-11.} Repository structure exists; \texttt{external/} verified
against the actual clone and \cref{sec:structure} corrected accordingly
(\cref{rem:treereality}); upstream pinned at \texttt{physicsnemo@59aaf59},
\texttt{physicsnemo-curator@86533e5}, \texttt{physicsnemo-cfd@0d2305e}, recorded in
\texttt{environment.md}. The compatibility table is populated from \emph{inspected data}
(\texttt{src/recon.py}, run on one case from each dataset) and the three dataset papers.
The canonical output space is committed: $[\Cp, \Cf_x, \Cf_y, \Cf_z]$ --- a choice that
is \emph{forced}, not preferred, because Windsor publishes no dimensional pressure
(\cref{rem:velocity}). The frontal-area convention is declared: \textbf{constant
reference throughout}, which means \texttt{force\_mom\_*.csv} for Ahmed and Windsor but
\texttt{force\_mom\_constref\_*.csv} for DrivAer (\cref{rem:refarea}).

All of the above is encoded with provenance in \texttt{src/datasets.py}, which is the
single source of truth for dataset differences. \textbf{No other file in this project may
hardcode a field name, a reference area, or a freestream velocity.}
\end{criterion}

\begin{remark}[What Phase 0 actually bought]
\label{rem:p0value}
Seven silent traps, none of which appeared in the plan as first written, and none of
which would have raised an error:

\begin{enumerate}[label=(\arabic*),leftmargin=2em]
  \item Variable ordering hardcoded at \texttt{loss.py:423} --- channel 0 must be
  pressure, channel 1 must be $x$-wall-shear (\cref{rem:ordering}).
  \item Fine-tuning via \texttt{resume\_dir} restores the epoch counter, so a
  fine-tuning run starting from a 500-epoch checkpoint with
  \texttt{train.epochs=50} exits immediately having trained nothing --- and writes a
  checkpoint. The recovery curve would have been a flat line (\cref{rem:epochtrap}).
  \item \texttt{force\_mom\_1.csv} means \emph{constant} reference in Ahmed and Windsor
  but \emph{variable} in DrivAer. 47 drag counts on Ahmed (\cref{rem:refarea}).
  \item Windsor ships \texttt{.vtu} with point data, not \texttt{.vtp} with cell data
  (\cref{rem:windsorfmt}).
  \item Windsor's wall shear is three scalars in \texttt{cfx}, \texttt{cfz},
  \texttt{cfy} order --- $z$ before $y$ (\cref{rem:windsoraxes}).
  \item $U_\infty$ spans 39$\times$; dimensional wall shear spans $\sim$1500$\times$
  (\cref{rem:velocity}).
  \item Windsor is \textbf{y-up} and \textbf{yawed $-2.5^\circ$}
  (\cref{rem:windsoraxes}).
\end{enumerate}

Every one of these produces a smoothly converging loss curve and a worthless model. Not
one produces an error message. This is the entire justification for the exit-criterion
discipline: \textbf{the dominant risks in this project are silent, and the criteria are
the only detectors.}

One trap was defused rather than found: the streamwise axis \emph{is} $x$ in all three
datasets (\cref{rem:streamwise}), so DoMINO's \texttt{normals[:, :, 0]} assumption holds.
That was verified, not assumed.
\end{remark}

% ============================================================
\section{Phase 1 --- Baseline Reproduction}
% ============================================================

\textbf{Goal:} demonstrate the pipeline reproduces the published DoMINO result on
DrivAerML. This is not a contribution --- it is the calibration of the instrument.
If the paper's numbers cannot be recovered, no later number means anything.

\subsection{Preprocess and compute scaling factors}

Use PhysicsNeMo-Curator (the supported path) to convert raw \texttt{stl}/\texttt{vtp}
into training-ready Zarr/NumPy. Then:

\begin{lstlisting}[style=sh]
python compute_statistics.py   # surface/volume field + coordinate statistics
\end{lstlisting}

This populates \texttt{surface\_factors}, which normalize labels to zero mean and
unit variance across the dataset. \textbf{Retain these; they become a leakage
hazard in \cref{sec:phase3}.}

\subsection{Train (surface only)}

\begin{lstlisting}[style=yml]
project: {name: domino_xgeom}
expt: p1_drivaer_baseline
data:
  input_dir:     /path/processed/drivaer/train
  input_dir_val: /path/processed/drivaer/val
  bounding_box_surface: [...]      # from compute_statistics.py
model:
  surface_points_sample: 8192      # reduce on OOM -- this is the knob
  geom_points_sample:    8192      # must be < n points on the STL
train: {epochs: 200}
\end{lstlisting}

\begin{lstlisting}[style=sh]
torchrun --nproc_per_node=8 train.py --config-name=config
\end{lstlisting}

NVIDIA reports full training reduced from roughly 5 days to just over 4 hours on
8$\times$H100. MSE loss is documented as giving the best results for both surface
and volume. The experimental matrix in \cref{sec:arms} is therefore affordable.

\subsection{Reproduction target}

Two published baselines exist and they differ. The \emph{paper} reports test-set
metrics over a split including OOD samples; the \emph{documentation} reports an
example validation run. Both are legitimate; they are different runs on different
splits. Record which you are targeting.

\begin{center}\small
\begin{tabular}{@{}lccc@{}}
\toprule
\textbf{Surface metric} & \textbf{Paper (rel.\ $L_2$)} & \textbf{Docs example run} & \textbf{Yours} \\
\midrule
Pressure         & 0.1505 & 0.101 & \\
X-wall-shear     & 0.2124 & 0.138 & \\
Y-wall-shear     & 0.3020 & 0.174 & \\
Z-wall-shear     & 0.3359 & 0.198 & \\
\midrule
\textbf{Drag $R^2$} & \textbf{0.96} & \textbf{0.983} & \\
Lift $R^2$       & --- & 0.971 & \\
\bottomrule
\end{tabular}
\end{center}

\begin{remark}[Do not average the two, and do not quietly pick the flattering one]
The gap between them (e.g.\ 0.1505 vs 0.101 on pressure) is real and reflects
different splits and training configurations. Target the \textbf{docs example run},
since it is the configuration shipped in the repository and therefore the one you are
actually reproducing --- and say so. If your number lands between the two, that is
fine; if it lands outside both, something is wrong.
\end{remark}

\begin{criterion}[Phase 1]
\label{crit:p1}
Surface pressure rel.\ $L_2$ within $\pm 20\%$ of the targeted baseline, and
\textbf{Drag $R^2 > 0.95$} on held-out DrivAer morphs. If not: stop and debug. Do not
proceed with a broken instrument.
\end{criterion}

% ============================================================
\section{Phase 2 --- Force Integration and Metrics Harness}
\label{sec:phase2}
% ============================================================

\textbf{Goal:} build the measurement apparatus once, correctly. Pure software; no
GPU.

\begin{remark}[This precedes the experiment, not follows it]
The temptation is to run the interesting transfer arm first and work out
measurement afterwards. If the force integrator carries a sign or area error, a bad
transfer number cannot be attributed to physics rather than arithmetic --- and the
distinction cannot be recovered after the fact.
\end{remark}

\subsection{Drag from a predicted surface field}

Drag is the streamwise component of integrated surface traction: pressure acting
along the outward normal, plus wall shear acting tangentially.
\begin{equation}
F_x = \underbrace{\oint_{\calS} \bigl(-p\,\hat{n}\bigr)\cdot\hat{x}\;dS}
        _{\text{pressure drag (dominant on a bluff body)}}
    \;+\;
      \underbrace{\oint_{\calS} \boldsymbol{\tau}_w\cdot\hat{x}\;dS}
        _{\text{viscous drag}}
\label{eq:forces}
\end{equation}

\begin{lstlisting}
"""forces.py -- surface integration. Single source of truth."""
import numpy as np, pyvista as pv

def surface_forces(surf, p_key, wss_key, rho, u_inf, frontal_area):
    surf  = surf.compute_normals(cell_normals=True, point_normals=False,
                                 auto_orient_normals=True)
    sized = surf.compute_cell_sizes(length=False, area=True, volume=False)
    areas   = sized.cell_data["Area"]        # (N,)
    normals = surf.cell_data["Normals"]      # (N,3) outward

    c   = surf.point_data_to_cell_data() if p_key in surf.point_data else surf
    p   = np.asarray(c.cell_data[p_key])     # (N,)
    wss = np.asarray(c.cell_data[wss_key])   # (N,3)

    q  = 0.5 * rho * u_inf**2
    Fp = (-(p[:, None] * normals) * areas[:, None]).sum(0)   # pressure
    Fv = ( wss                    * areas[:, None]).sum(0)   # viscous
    F  = Fp + Fv
    den = q * frontal_area
    return {"Cd": float(F[0]/den), "Cd_pressure": float(Fp[0]/den),
            "Cd_viscous": float(Fv[0]/den), "Cl": float(F[2]/den)}
\end{lstlisting}

\begin{remark}[Sign conventions will bite]
Normal orientation (inward vs outward), the sign of the pressure term, and which
axis is streamwise vary between datasets and between VTK writers. Validate against
ground truth before trusting a single prediction. This is not optional.
\end{remark}

\subsection{Validating the integrator}

Each dataset ships time-averaged force coefficients. Run the integrator on the
\emph{ground-truth} surface fields, for $\geq 20$ cases per dataset, and compare.

\begin{center}\small
\begin{tabular}{@{}p{2.4cm}p{3.2cm}p{9cm}@{}}
\toprule
\textbf{Result} & \textbf{Relative $\Cd$ error} & \textbf{Action} \\
\midrule
Pass         & $< 2\%$   & Proceed. \\
Investigate  & $2$--$5\%$ & Likely cell-area or normal-orientation issue. \\
Fail         & $> 5\%$   & Sign error or wrong reference area. \textbf{Do not proceed.} \\
\bottomrule
\end{tabular}
\end{center}

Cross-check against PhysicsNeMo-CFD, which contains the force-computation utilities
used to produce the published drag-$R^2$ plots. Where the two disagree, PhysicsNeMo-CFD
is correct.

\subsection{Metrics}

\begin{lstlisting}
"""metrics.py -- the only three numbers that matter."""
import numpy as np
from scipy.stats import spearmanr

def drag_counts(cd_pred, cd_true):        # 1 count = 0.001 Cd
    return float(np.mean(np.abs(np.asarray(cd_pred)-np.asarray(cd_true)))*1000)

def design_ranking(cd_pred, cd_true):     # Definition 3.2 -- usability
    rho, p = spearmanr(cd_pred, cd_true)
    return {"spearman_rho": float(rho), "p_value": float(p)}

def cp_field_error(cp_pred, cp_true, areas):   # Definition 3.3 -- honesty
    w = areas / areas.sum(); e = cp_pred - cp_true
    return {"l2_area_weighted": float(np.sqrt(np.sum(w*e**2))),
            "mean_signed_bias": float(np.sum(w*e))}   # <- cancellation detector
\end{lstlisting}

\begin{criterion}[Phase 2]
\label{crit:p2}
The integrator reproduces published $\Cd$ from ground-truth fields to within $2\%$
on all three datasets. Committed and unit-tested.
\end{criterion}

% ============================================================
\section{Phase 3 --- Unified Multi-Body Datapipe}
\label{sec:phase3}
% ============================================================

\textbf{Goal:} make Ahmed, Windsor, and DrivAer indistinguishable to DoMINO. This
is the hardest engineering in the project and where silent bugs live.

\subsection{Per-dataset path classes}

The ``Extending DoMINO to a custom dataset'' documentation is the map. Modify
\texttt{openfoam\_datapipe.py}, following the pattern NVIDIA gives:

\begin{lstlisting}
# One class per body. Fill from the Phase-0 recon output -- never from memory.
class AhmedPaths:
    @staticmethod
    def geometry_path(case_dir): return case_dir / "<VERIFY>.stl"
    @staticmethod
    def surface_path(case_dir):  return case_dir / "<VERIFY>.vtp"

class WindsorPaths:
    @staticmethod
    def geometry_path(case_dir): return case_dir / "<VERIFY>.stl"
    @staticmethod
    def surface_path(case_dir):  return case_dir / "<VERIFY>.vtp"

class DrivAerPaths:      # shipped in-repo; use as reference implementation
    ...
\end{lstlisting}

Enforce \cref{eq:canonical} at the datapipe boundary --- see \cref{rem:ordering}.
Assert it; do not trust it.

\subsection{The bounding-box decision}

DoMINO scales input coordinates by user-defined bounding boxes
(\texttt{data.bounding\_box\_surface}), with \texttt{normalize\_coordinates}
enabled by default; the presets are tuned for DrivAerML. Two options, and this is
a \emph{research} decision, not a config detail:

\begin{center}\small
\begin{tabular}{@{}p{2.2cm}p{5.4cm}p{6.4cm}@{}}
\toprule
\textbf{Option} & \textbf{Mechanism} & \textbf{Consequence} \\
\midrule
A (\textbf{chosen}) & Per-body bounding boxes; each body normalized into its own box
  & Model sees all bodies at unit scale. Makes shape transfer the object of study;
    discards absolute size. \\
B & One global box enclosing the largest body
  & Preserves relative scale, but Ahmed occupies a small corner of DrivAer's box;
    the model may never learn to use the full domain. \\
\bottomrule
\end{tabular}
\end{center}

\textbf{Option A is selected} because the question is about \emph{shape}
generalization, not scale generalization. Document the choice; when asked ``did you
normalize scale away?'', the answer must be a confident yes, with a reason.

\subsection{Scaling factors and target leakage}

\begin{remark}[Target-statistics leakage --- subtle enough to survive review]
\label{rem:leak}
\texttt{compute\_statistics.py} produces \texttt{surface\_factors} over ``the''
dataset. For a mixed Ahmed+Windsor training set, compute them over the
\textbf{combined source set only}, and apply those \emph{same, unchanged} factors to
DrivAer at test time. Recomputing statistics on DrivAer leaks target information
into a supposedly cold transfer. The resulting number would be optimistic, would
look entirely reasonable, and would invalidate the central claim of the study.

\textbf{This applies per source set.} Arm A3a trains on Ahmed alone and therefore has
its \emph{own} factors, computed over Ahmed only. A3 has factors computed over
Ahmed\,+\,Windsor. They are different files and must not be interchanged; using A3's
factors for A3a would leak Windsor statistics into the Ahmed-only arm and destroy the
comparison that \cref{hyp:diversity} rests on.
\end{remark}

\begin{criterion}[Phase 3]
\label{crit:p3}
A single assertion passes on a batch drawn from each of the three bodies: field
names, variable ordering, units, and normalization are identical. Written as a test,
not as a belief.
\end{criterion}

% ============================================================
\section{Phase 4 --- The Experiment}
\label{sec:arms}
% ============================================================

\subsection{The arms}

\begin{center}\small
\begin{tabular}{@{}p{1.1cm}p{3.4cm}p{3.0cm}p{7.0cm}@{}}
\toprule
\textbf{Arm} & \textbf{Train} & \textbf{Test} & \textbf{What it establishes} \\
\midrule
\textbf{A1} & DrivAer (400 morphs) & DrivAer (100 held out)
  & Baseline. The instrument works. Interpolation only. \\[0.3em]
\textbf{A2} & DrivAer (all 500) & Ahmed + Windsor
  & \emph{Downward} (\cref{hyp:downward}). Control; expect success.
    \textbf{Gate: run before A3.} \\[0.3em]
\textbf{A3a} & Ahmed only (500) & DrivAer (500)
  & Interpretive control (\cref{hyp:diversity}). Does adding a second body family
    buy anything? \\[0.3em]
\textbf{A3} & \textbf{Ahmed + Windsor (855)} & \textbf{DrivAer (500)}
  & \textbf{The experiment} (\cref{hyp:degrade}). Upward transfer. \\
\bottomrule
\end{tabular}
\end{center}

\begin{remark}[Why A3a exists, and what it is not]
\label{rem:a3a}
With A3 alone, a failure is ambiguous: did transfer break because the \emph{topology}
is unseen, or because two source bodies is not a distribution? A3a resolves that
ambiguity. If A3 materially beats A3a, source diversity is the currency and the study's
sentence becomes ``geometric diversity buys transfer, at roughly this rate.'' If the two
are indistinguishable, adding bluff-body variety does not help, and the barrier is
topological --- a cleaner and stronger reading of \cref{hyp:degrade}.

\textbf{Two points are not a curve.} A3a licenses no extrapolation to ``how many body
families are needed.'' Do not draw a line through it.

\textbf{The confound is accepted, not closed.} Ahmed alone is 500 cases; Ahmed +
Windsor is 855. If A3 beats A3a, diversity and sample count are not separated. Closing
this would require a third arm (Ahmed + Windsor subsampled to 500, matched count,
higher diversity); it is \emph{deliberately not run}, because it is the first step of
the diversity-scaling study, which is a different project (\cref{rem:scope}). State the
confound in the limitations rather than spending a run on it.
\end{remark}

\begin{remark}[A2 is not filler --- it is what makes A3 believable]
If only A3 were run and it failed, a skeptic says: your pipeline is broken. A2
succeeding \emph{on the same pipeline} forecloses that objection. It is a control,
and it is the reason the negative result will be believed. Run it \emph{before} A3:
if the downward control fails, A3 is uninterpretable and there is no point spending
the run.
\end{remark}

\begin{remark}[The fine-tuning epoch trap --- would silently flatten the recovery curve]
\label{rem:epochtrap}
There is no \texttt{retraining.py}. Fine-tuning is done by pointing \texttt{resume\_dir}
at a pre-trained checkpoint and running \texttt{train.py}. Doing this \emph{naively}
produces a silently wrong result, and the recovery curve --- the study's headline figure
--- is what it corrupts.

In \texttt{train.py} (line 646 at \texttt{59aaf59}):
\begin{lstlisting}
init_epoch = load_checkpoint(
    to_absolute_path(cfg.resume_dir),
    models=model, optimizer=optimizer,
    scheduler=scheduler, scaler=scaler, device=dist.device,
)
if init_epoch != 0:
    init_epoch += 1
epoch_number = init_epoch
\end{lstlisting}

\textbf{Two failures follow.} (i) The A3 checkpoint was saved at epoch 500. Setting
\texttt{train.epochs=50} for a fine-tuning run starts the loop at 501, which exceeds 50,
so it \textbf{exits immediately} --- having trained for zero steps --- and writes a
checkpoint. No crash, no warning. Every fine-tuned model at every $N$ would be identical
to the cold model, the recovery curve would be a flat line, and the conclusion would be
``fine-tuning does not help.'' (ii) The optimizer and scheduler are restored as well, so
``fine-tune at a small learning rate'' is not something you get for free --- you inherit
whatever the schedule had decayed to.

\textbf{The fix is small, because \texttt{load\_checkpoint} skips arguments that are
\texttt{None}} (\texttt{physicsnemo/utils/checkpoint.py}, line 964: \emph{``Objects that
are \texttt{None} are silently skipped''}). Patch \texttt{train.py} so that, under a
fine-tuning flag, it (a) calls \texttt{load\_checkpoint} with \texttt{models=model} only
--- no optimizer, no scheduler, no scaler --- and (b) \textbf{forces
\texttt{init\_epoch = 0}}, discarding the returned value. The fine-tuning learning rate
then comes from the config, as intended. Gate it behind a flag so the ordinary
crash-resume path is unchanged.

This is \textbf{patch \#1} in \texttt{patches/}. \Cref{rem:patches} was right that the
delta would be small; it was wrong that it would be confined to the datapipe.
\end{remark}

\subsection{The recovery curve}

From the A3 checkpoint, fine-tune on $N$ DrivAer cases, $N \in \{0,5,10,25,50,100\}$,
evaluating throughout on a \textbf{fixed DrivAer holdout that is never fine-tuned on}
(100 cases, reserved from the outset). The recovery curve is built from A3 only; A3a
contributes a single cold ($N{=}0$) number, for comparison against A3's cold number.

The repository already implements this. The documented retraining recipe: place the
pretrained checkpoint in \texttt{resume\_dir}, place the surface scaling factors in
the output directory, run \texttt{retraining.py} for a specified number of epochs at
a small learning rate from the checkpoint, then \texttt{test.py} --- which runs on
\emph{raw} \texttt{.vtp}/\texttt{.vtu} and writes predictions back into the file.

\begin{lstlisting}[style=sh]
# 1. pretrained A3 checkpoint -> resume_dir  (conf/config.yaml)
# 2. SOURCE-set scaling factors -> output dir   (see Remark 8.1: do not recompute)
python retraining.py ++train.epochs=<E> ++data.input_dir=<N-case subset>
python test.py       # inference on raw vtp/vtu; predictions written in-place
\end{lstlisting}

\begin{remark}[Three seeds per $N$, or the curve is an anecdote]
Run each $N$ with three independent random draws of \emph{which} $N$ cases are used.
At $N=5$, the identity of the five cases will dominate the outcome. A single-seed
curve is not evidence. Plot mean and spread.
\end{remark}

\subsection{Logging schema}

\begin{lstlisting}
{"arm": "A3",     # one of A1 | A2 | A3a | A3
 "n_finetune": 25, "seed": 1,
 "cd_mae_counts":      ...,   # headline
 "cd_pressure_counts": ...,   # decomposition: where does the error live?
 "cd_viscous_counts":  ...,
 "spearman_rho":       ...,   # usability   (Def. 3.2)
 "cp_l2_area_wtd":     ...,   # honesty     (Def. 3.3)
 "cp_signed_bias":     ...,   # cancellation detector (Remark 3.2)
 "drag_r2":            ...}
\end{lstlisting}

\begin{criterion}[Phase 4]
\label{crit:p4}
Arms run \emph{in order}: A1, then A2 (gate --- if the downward control fails, the
pipeline is broken and A3 is uninterpretable; stop), then A3a, then A3. Recovery curve
populated at 6 values of $N$ $\times$ 3 seeds from the A3 checkpoint. Every run logged
to the schema above. \cref{hyp:diversity} adjudicated from the A3a-vs-A3 cold numbers.
\end{criterion}

% ============================================================
\section{Phase 5 --- Error Localization}
% ============================================================

\textbf{Goal:} establish \emph{why} it fails, not merely \emph{that} it fails. This
is the section that distinguishes an engineer from someone who ran a training script,
and it is where \cref{hyp:localize} is adjudicated.

\subsection{Surface zones}

Segment the DrivAer surface once, by hand, and report error per zone:

\begin{center}\small
\begin{tabular}{@{}p{4.6cm}p{3.0cm}p{7.6cm}@{}}
\toprule
\textbf{Zone} & \textbf{Present in source?} & \textbf{Expectation under \cref{hyp:localize}} \\
\midrule
Front end / stagnation      & Yes (both) & Accurate \\
Hood, windscreen, roof      & Yes (both) & Accurate \\
\textbf{C-pillar / rear window} & \textbf{No} & \textbf{Large error} (notchback separation) \\
\textbf{Rear end / base wake}   & Partly (Windsor) & Moderate error \\
\textbf{Wheels, wheelhouses}    & \textbf{No} & \textbf{Large error} \\
\textbf{Mirrors}                & \textbf{No} & \textbf{Large error} \\
Underbody                   & No & Large error \\
\bottomrule
\end{tabular}
\end{center}

\begin{remark}[The prediction is registered in advance]
\Cref{hyp:localize} is stated before the analysis is run. If it holds, the study
has a \emph{mechanistic} explanation rather than a number. If it does not hold, that
is the more interesting outcome, and it was found honestly rather than rationalized
afterwards. Do not revise the hypothesis after seeing the error map.
\end{remark}

\subsection{Wake evidence: showing the mechanism, not asserting it}
\label{sec:wake}

\begin{remark}[Why surface-only is not sufficient \emph{here}, and only here]
\label{rem:wakegap}
This study is surface-only by design (\cref{rem:novtu}) --- drag requires only surface
integration, and the volume data is a factor of $\sim$1000 larger. That decision is
correct for the \emph{metric}. It is not sufficient for the \emph{mechanism}.

\Cref{hyp:localize} and the argument of \cref{sec:whybreak} claim the model fails
because it has no learned representation of a C-pillar vortex or a wheel wake. With
surface data alone you can show \emph{that} the C-pillar pressure is wrong. You cannot
show \emph{why} --- because the cause lives in the flow \emph{off} the surface. The
reviewer's question writes itself: \emph{``You claim it fails to represent the C-pillar
vortex. Show me the vortex.''} Without volume data for at least a few cases, there is
no answer.
\end{remark}

\textbf{The bounded fix.} Fetch VTU volume data for \textbf{five DrivAer cases only}:

\begin{center}\small
\begin{tabular}{@{}p{2.2cm}p{4.4cm}p{8.6cm}@{}}
\toprule
\textbf{Cases} & \textbf{Selection} & \textbf{Purpose} \\
\midrule
2 & Worst cold-transfer $\Cd$ error & Where the mechanism should be most visible \\
2 & Best cold-transfer $\Cd$ error  & Control: is the wake right when the drag is right? \\
1 & Median error                    & Reference \\
\bottomrule
\end{tabular}
\end{center}

Compare predicted against true velocity fields, \emph{qualitatively}, in three places:
the C-pillar / rear-window separation region, the wheelarch, and the base wake. Two
planes suffice (a centreline $x$--$z$ cut and an $x$--$y$ cut through the wheel).

\begin{remark}[Scope guard --- this is evidence, not a second study]
Five cases. Qualitative comparison. No volume error metrics in the headline, no volume
model trained, no volume term in the loss. The purpose is a single figure panel that
makes \cref{hyp:localize} \emph{shown} rather than \emph{asserted}. If this begins to
grow into a volume-field study, stop: that is a different project
(\cref{rem:scope}). Storage cost is negligible at five cases; at five hundred it is the
failure mode \cref{rem:novtu} exists to prevent.
\end{remark}

\subsection{The three figures}

There are only three. Anything beyond them is dilution.

\begin{enumerate}[label=F\arabic*.]
  \item \textbf{The Recovery Curve.} $x$: fine-tuning budget $N$. $y$: $\Cd$ MAE in
  drag counts. Lines: A1 baseline (horizontal reference), A3 cold at $N{=}0$, A3
  fine-tuned. Shaded band for seed spread. \textbf{Annotate the knee.} Mark A3a's cold
  number as a single point at $N{=}0$ alongside A3's, so the diversity comparison
  (\cref{hyp:diversity}) is visible on the money figure --- but do \emph{not} connect
  them with a line (\cref{rem:a3a}). This is the money figure --- it carries the
  deliverable sentence.

  \item \textbf{The Honesty Check.} Surface $\Cp$ error painted on the DrivAer body;
  three views (side, rear, underside); diverging colormap centred at zero so sign is
  visible. Two panels: cold ($N{=}0$) vs fine-tuned at the knee. The reader should
  \emph{see} error concentrate on wheels, mirrors, and C-pillar, then recede.

  \textbf{Plus one wake panel} (\cref{sec:wake}): for the worst cold-transfer case,
  predicted vs true velocity on a centreline cut through the C-pillar separation. This
  is what converts \cref{hyp:localize} from a claim into evidence --- the surface map
  shows \emph{where} the pressure is wrong; the wake cut shows \emph{why}.

  \item \textbf{The Ranking Test.} Predicted vs true $\Cd$ across held-out designs;
  diagonal reference; Spearman $\rho$ inset. Two panels, cold vs fine-tuned. This is
  where a reader decides whether they could actually optimize with this model.
\end{enumerate}

\begin{criterion}[Phase 5]
\label{crit:p5}
Three figures exist, F2 including the wake panel; the per-zone error table is
populated; \cref{hyp:localize} is explicitly adjudicated as held or refuted; and the
failure location \emph{and its physical cause} can be stated in one sentence ---
with the wake cut as the evidence for the second half of that sentence.
\end{criterion}

% ============================================================
\section{Phase 6 --- Deliverable}
% ============================================================

\begin{remark}[Design constraint]
The reader will spend approximately eight minutes on this artifact and will not read
a methods section. The README \emph{is} the deliverable; the code is the evidence.
Structure accordingly.
\end{remark}

\subsection{Structure}

\begin{enumerate}[label=\arabic*.]
  \item The question. One sentence.
  \item The answer. One sentence, with the number in it.
  \item Figure F1 (recovery curve).
  \item \textbf{Why this should be believed}: baseline reproduction
  (\cref{crit:p1}) and force-integrator validation (\cref{crit:p2}). Two rows each.
  Credibility is a down-payment; spend it early.
  \item Figures F2, F3.
  \item What it means for practice: the $N$-runs recommendation.
  \item Limitations, stated before anyone else finds them.
  \item Reproduction instructions.
\end{enumerate}

\subsection{The target sentence}

\begin{quote}
\emph{A DoMINO surrogate trained only on simple bluff bodies predicts DrivAer drag
to within \underline{\hspace{1.2cm}} counts and preserves design ranking at Spearman
$\rho =$ \underline{\hspace{1cm}} --- but only after fine-tuning on
\underline{\hspace{0.9cm}} high-fidelity runs. Below that, it obtains the right drag
for the wrong reasons: errors concentrate on the wheels and C-pillar and cancel in
the integral.}
\end{quote}

The blanks are unknown. That is the point, and it is what makes the study worth
running.

\subsection{Limitations to state yourself}

\begin{itemize}
  \item Surface-only for all metrics; volume data used qualitatively, on five cases,
  for mechanism only (\cref{sec:wake}).
  \item A single architecture; no claim that other surrogates behave the same way.
  \item Two source bodies is a thin ``distribution'' of geometries. A3a
  (\cref{rem:a3a}) tests whether the second body buys anything, but two points do not
  establish how much geometric diversity transfer requires --- and because Ahmed alone
  (500) and Ahmed\,+\,Windsor (855) differ in sample count as well as diversity, any
  gap between A3a and A3 does not separate the two causes. The diversity-scaling
  question is left open, deliberately.
  \item Geometry variation only, at fixed inlet conditions --- a limitation the
  AhmedML authors flag about their own dataset.
  \item No uncertainty quantification: the model does not know when it is
  extrapolating. This is the natural follow-up study.
\end{itemize}

\subsection{Declared future work: where does the failure live?}
\label{sec:ablation}

\begin{remark}[The most interesting experiment on the table --- and why it is not in this study]
\label{rem:ablation}
Once A3 exists, one ablation is cheap and sharper than anything else available: freeze
the \textbf{geometry encoder} and fine-tune only the \textbf{aggregation network}, then
the reverse. This separates two very different failures.

\emph{If freezing the encoder and tuning aggregation recovers performance:} the geometry
representation was adequate --- the encoder \emph{did} encode the wheel --- and only the
mapping from local encoding to surface pressure was wrong. Geometric generalization is
fine; the flow mapping is what does not transfer.

\emph{If it does not recover, and the reverse does:} the encoder itself is impoverished.
The learned point-convolution kernels (\cref{eq:pointconv}) genuinely failed to
represent an unseen feature. This is the strong form of \cref{hyp:degrade}, and it
would be a statement about neural operators in general, not about drag.

\textbf{It is deliberately excluded.} Not because it is uninteresting --- because it is
\emph{too} interesting. Its headline (\emph{``where does geometric generalization break
in a neural operator?''}) is a different and arguably better paper than this one
(\emph{``can it predict drag on an unseen car, and what does recovery cost?''}).
Pursuing both would give a reader two half-arguments instead of one whole one, and this
artifact has roughly eight minutes to land a single sentence (\cref{rem:scope}).

Run it \emph{after} the study ships, as the follow-up it deserves to be. Naming it here
--- rather than omitting it and hoping nobody asks --- is the difference between scope
discipline and an oversight.

\textbf{The same applies to diversity scaling.} A3a (\cref{rem:a3a}) gives two points on
a curve that would need many more to be a curve. ``How many source body families does
geometric transfer require, and does sample count substitute for diversity?'' is the
other follow-up. It is named here for the same reason.
\end{remark}

% ============================================================
\section{Risk Register}
% ============================================================

\begin{center}\small
\begin{tabular}{@{}p{4.0cm}p{1.4cm}p{9.6cm}@{}}
\toprule
\textbf{Risk} & \textbf{Phase} & \textbf{Mitigation} \\
\midrule
Variable ordering differs across datasets; loss trains against wrong targets
  & 0, 3
  & \cref{rem:ordering}. Verify per dataset; enforce \cref{eq:canonical}; assert in
    the datapipe. \textbf{Failure is silent --- loss curve looks healthy.} \\[0.3em]
Sign / normal-orientation error in force integration
  & 2
  & Validate against published GT forces to $<2\%$ \emph{before} trusting any
    prediction. Cross-check PhysicsNeMo-CFD. \\[0.3em]
Target-statistics leakage via scaling factors
  & 3
  & \cref{rem:leak}. Compute \texttt{surface\_factors} on the source set only; apply
    unchanged to DrivAer. Separate factors per source set (A3a vs A3). \\[0.3em]
A3 fails so completely the cold number is uninterpretable
  & 4
  & \textbf{Anticipated} (\cref{hyp:degrade}). The recovery curve is the deliverable;
    build it from the start, not as a rescue. \\[0.3em]
A3a$\approx$A3 read as ``diversity does not help'' when it is a sample-count artefact
  & 4
  & \cref{rem:a3a}. The confound is accepted and stated, not silently resolved. Do not
    claim a diversity effect the design cannot separate. \\[0.3em]
Recovery curve dominated by seed at small $N$
  & 4
  & Three seeds per $N$; report spread, never a single line. \\[0.3em]
Storage blowup ($\sim$160M-cell DrivAerML volumes)
  & 0
  & Surface-only. Download STL + VTP; skip VTU. \\[0.3em]
Scope creep (fidelity transfer, UQ, volume fields, diversity scaling)
  & All
  & \cref{rem:scope}. Explicitly deferred. \\[0.3em]
OOM during training
  & 1, 4
  & Reduce \texttt{model.surface\_points\_sample} --- the documented knob. \\
\bottomrule
\end{tabular}
\end{center}

% ============================================================
\section{Common Pitfalls}
% ============================================================

\begin{center}\small
\begin{tabular}{@{}p{4.2cm}p{4.0cm}p{6.8cm}@{}}
\toprule
Pitfall & Symptom & Fix \\
\midrule
Wrong variable order & Loss converges; model worthless
  & Enforce $[\Cp, \Cf_x, \Cf_y, \Cf_z]$ at datapipe boundary \\
Normals inward not outward & $\Cd$ sign flipped
  & \texttt{auto\_orient\_normals=True}; validate against GT \\
Fields on points, not cells & Area weighting wrong
  & \texttt{point\_data\_to\_cell\_data()} before integrating \\
Wrong reference area & $\Cd$ off by constant factor
  & Declare $A_{\text{frontal}}$ convention; match dataset authors' \\
Scaling factors recomputed on target & Cold transfer looks too good
  & Source-set factors only (\cref{rem:leak}) \\
A3's factors reused for A3a & Windsor stats leak into Ahmed-only arm
  & Separate factor files per source set (\cref{rem:leak}) \\
Right $\Cd$, wrong $\Cp$ field & $\mathcal{E}_{L2}$ large, $\mathcal{E}_{bias}\approx 0$
  & Not a bug --- this is \cref{hyp:cancel}. Report it; it is a finding \\
Good $\Cd$ MAE, poor $\rho$ & Model unusable for optimization
  & Report both; $\rho$ is the deployability gate (\cref{def:rank}) \\
Single-seed recovery curve & Non-monotone, noisy
  & 3 seeds per $N$ \\
Line drawn through A3a and A3 & Implies a diversity scaling law
  & Two points. Plot as points, not a line (\cref{rem:a3a}) \\
OOM on DrivAer, not on Ahmed & Crash mid-arm
  & Lower \texttt{surface\_points\_sample} for \emph{all} arms, not just one \\
Volume data downloaded & Tens of TB consumed
  & Surface-only project; STL + VTP suffice \\
\bottomrule
\end{tabular}
\end{center}

% ============================================================
\section{Quick Reference}
% ============================================================

\textbf{Question:} Ahmed + Windsor $\to$ DrivAer. Can it predict drag on a body it
has never seen, and what does repair cost?

\textbf{Architecture (3 stages):} (1) \emph{Global} --- STL point cloud $\to$ structured
grid via learnable ball-query point convolutions $y_i = \sum_j f(\vec{x}_i, \vec{x}_j,
d_{ij})$ at multiple radii; propagated to the domain grid by iterative CNN blocks;
SDF + gradients appended. (2) \emph{Local} --- per query point, extract an
$l\times l\times l\times n_f$ sub-region encoding from the global encoding. (3)
\emph{Aggregate} --- basis-function net over a $p+1$ stencil (coords, SDF, normals)
$\to$ concat with local encoding $\to$ inverse-distance-weighted solution.
Loss: per-variable MSE + area-weighted MSE. Adam, 500 epochs, AMP fp16.

\textbf{Why it should break:} the kernels $f$ are \emph{learned}. No wheel, mirror, or
C-pillar ever passes through $f$ during Ahmed/Windsor training
(\cref{sec:whybreak}).

\textbf{Not the same as the paper's OOD:} theirs $=$ drag outside training range, same
body. Ours $=$ different body (\cref{rem:ood}). Say this explicitly.

\textbf{Canonical target:} $\mathbf{y} = [\Cp, \Cf_x, \Cf_y, \Cf_z]$, fixed order,
per-body nondimensionalization (\cref{eq:canonical}).

\textbf{Metrics:} $\Cd$ MAE in drag counts (accuracy); Spearman $\rho$ (usability);
$\mathcal{E}_{L2}$ and $\mathcal{E}_{bias}$ (honesty, \cref{eq:honesty}).

\textbf{Arms (in order):} A1 DrivAer$\to$DrivAer (instrument check); A2
DrivAer$\to$Ahmed+Windsor (control, \emph{gate}); A3a Ahmed$\to$DrivAer (interpretive
control, \cref{rem:a3a}); A3 Ahmed+Windsor$\to$DrivAer (\emph{the experiment}).

\textbf{Recovery:} fine-tune A3 on $N \in \{0,5,10,25,50,100\}$, 3 seeds each, fixed
never-fine-tuned holdout.

\textbf{Figures:} F1 recovery curve (annotate the knee; A3a as a bare point at $N{=}0$,
\emph{not} joined to A3); F2 $\Cp$ error map (cold vs fine-tuned); F3 ranking scatter
with $\rho$.

\textbf{Hypothesis status:} \cref{hyp:degrade} registered, expected to hold;
\cref{hyp:localize} registered, adjudicate in Phase 5; \cref{hyp:cancel} registered,
adjudicate via $\mathcal{E}_{bias}$; \cref{hyp:downward} registered as control;
\cref{hyp:diversity} registered as interpretive control (A3a vs A3).
\textbf{None are proved. None may be cited as results.}

\textbf{Exit criteria:} \cref{crit:p0} compatibility table; \cref{crit:p1} Drag
$R^2 > 0.95$; \cref{crit:p2} integrator within $2\%$; \cref{crit:p3} cross-body
assertion passes; \cref{crit:p4} all arms + curve; \cref{crit:p5} three figures +
hypothesis adjudicated.

% ============================================================
\appendix
\section{Project Structure}
\label{sec:structure}

\begin{remark}[Epistemic status of this section]
\label{rem:structstatus}
This is a \textbf{proposal to be created in Phase 0}, not a description of something
that exists. The layout of \emph{your} repository is a design choice, argued for below.
The internal layout of \texttt{external/physicsnemo/.../domino/} is
\textbf{taken from documentation, not from having read the tree} --- the documentation
names \texttt{conf/config.yaml}, \texttt{conf/cached.yaml}, \texttt{src/conf},
\texttt{train.py}, \texttt{test.py}, \texttt{retraining.py},
\texttt{compute\_statistics.py}, \texttt{cache\_data.py}, \texttt{process\_data.py},
\texttt{openfoam\_datapipe.py}, \texttt{inference\_on\_stl.py}, and
\texttt{src/shuffle\_volumetric\_curator\_output.py}, but the exact directory nesting
is unverified. \textbf{Verify on clone and correct this section.} Consistent with the
discipline of this document: what is known is stated, what is assumed is labelled.
\end{remark}

\subsection{The two-repository principle}

The single most important structural decision is that \textbf{their code and your code
do not mix}. PhysicsNeMo is cloned as a pinned dependency, never forked, never edited
in place. Your repository contains configuration, analysis, tests, and a small,
explicit set of patches.

\begin{remark}[Why this matters, concretely]
\label{rem:twolevel}
Three reasons, in ascending order of importance.

\emph{(i) The upstream moves.} PhysicsNeMo is under active development; the
documentation itself notes performance work reducing training from $\sim$5 days to
$\sim$4 hours across releases, and that physics-loss support ``might introduce
breaking changes.'' If your changes are smeared across a fork, re-basing onto a new
release is archaeology.

\emph{(ii) Your contribution becomes invisible.} If you edit
\texttt{openfoam\_datapipe.py} in place inside a clone, your entire engineering
contribution --- which is genuinely small, and that is a \emph{virtue} --- disappears
into a diff against a 100k-line repository. A reviewer, an interviewer, or you in six
months cannot see what you did.

\emph{(iii) Reproducibility.} ``Clone PhysicsNeMo at commit \texttt{abc123}, apply
these two patches, run these configs'' is a reproducible instruction. ``Here is my
fork'' is not.
\end{remark}

\subsection{Your repository}

\begin{lstlisting}[style=sh]
~/domino-xgeom/                     # git init THIS. Small. Auditable.
|
|-- README.md                       # the eight-minute deliverable (Phase 6)
|-- plan.tex  plan.pdf              # THIS document, kept current (see 15.6)
|-- environment.md                  # commit hashes, CUDA, GPU type, driver
|-- .gitignore                      # external/, data paths, *.pt, *.zarr
|
|-- external/                       # THEIRS. Pinned. Gitignored. Never edited.
|   |-- physicsnemo/                #   model, train.py, test.py, retraining.py
|   |-- physicsnemo-curator/        #   preprocessing -> zarr/npy
|   `-- physicsnemo-cfd/            #   force utils; Phase-2 cross-check oracle
|
|-- patches/                        # ** YOUR ONLY EDITS TO THEIR CODE **
|   |-- openfoam_datapipe.py        #   + AhmedPaths, WindsorPaths
|   |-- train_losses.py             #   IF variable order needs fixing (Rem 4.1)
|   |-- apply.sh                    #   copy patches -> external/; idempotent
|   `-- UPSTREAM.md                 #   what was changed, why, against which commit
|
|-- conf/                           # YOUR Hydra configs. One per arm. Diffable.
|   |-- base.yaml                   #   shared: bounding boxes, sampling, epochs
|   |-- a1_drivaer_baseline.yaml    #   Arm A1  (instrument check)
|   |-- a2_downward.yaml            #   Arm A2  (control; GATE -- run before A3)
|   |-- a3a_ahmed_only.yaml         #   Arm A3a (interpretive control, Rem 11.1)
|   |-- a3_upward.yaml              #   Arm A3  (** the experiment **)
|   `-- ft_N{0,5,10,25,50,100}.yaml #   recovery curve; seed passed at runtime
|
|-- src/                            # YOUR analysis code. None of this is DoMINO.
|   |-- recon.py                    #   Phase 0: dataset compatibility
|   |-- forces.py                   #   Phase 2: SINGLE SOURCE OF TRUTH for Cd
|   |-- metrics.py                  #   Phase 2: drag counts, Spearman, signed bias
|   |-- zones.py                    #   Phase 5: surface segmentation
|   |-- collect.py                  #   Phase 4: results/*.json -> one dataframe
|   `-- figures.py                  #   Phase 5: F1, F2, F3
|
|-- tests/                          # Exit criteria, executable.
|   |-- test_forces.py              #   GT integration within 2%   -> Crit 8.1
|   `-- test_datapipe.py            #   cross-body assertion       -> Crit 9.1
|
|-- results/                        # APPEND-ONLY. One JSON per run. 22 files.
|   `-- {arm}_{N}_{seed}.json       #   schema: Section 11.4
|
`-- figures/
    |-- F1_recovery_curve.pdf
    |-- F2_cp_error_map.pdf
    `-- F3_ranking_scatter.pdf
\end{lstlisting}

\subsection{Storage (outside the repository)}

Datasets and checkpoints do not belong in git. They live on scratch/fast storage:

\begin{lstlisting}[style=sh]
/scratch/domino-xgeom/
|
|-- raw/                            # STL + VTP ONLY.  NO VTU.  (see Rem 15.3)
|   |-- ahmed/     <case>/{*.stl, *.vtp}
|   |-- windsor/   <case>/{*.stl, *.vtp}
|   `-- drivaer/   <case>/{*.stl, *.vtp}
|
|-- wake/                           # ** EXCEPTION: 5 cases only. Phase 5. **
|   `-- drivaer/   <case>/*.vtu     #    2 worst + 2 best + 1 median (Sec 12.2)
|                                   #    fetched AFTER A3, once errors are known
|
|-- processed/                      # physicsnemo-curator output
|   |-- source_train/               #   Ahmed + Windsor, mixed   -> A3
|   |-- source_val/
|   |-- ahmed_only_train/           #   Ahmed alone              -> A3a
|   |-- ahmed_only_val/
|   |-- drivaer_train/              #   -> A1
|   |-- drivaer_val/
|   |-- drivaer_holdout/            # ** NEVER fine-tuned on. Never. **
|   `-- ft_pools/N{5,10,25,50,100}_seed{1,2,3}/
|
|-- scaling/
|   |-- source_factors.json         # ** Ahmed+Windsor ONLY. See Rem 15.4 **
|   |-- ahmed_only_factors.json     # ** Ahmed ONLY. A3a. Do NOT reuse above. **
|   `-- drivaer_factors.json        #    for Arm A1 baseline only
|
`-- checkpoints/
    |-- a1/  a2/  a3a/  a3/
    `-- ft_N{...}_seed{...}/
\end{lstlisting}

\subsection{The structural decisions that are not cosmetic}

\begin{remark}[\texttt{patches/} keeps your contribution visible and re-appliable]
\label{rem:patches}
Your entire code modification to DoMINO is two files: path classes for Ahmed and
Windsor in \texttt{openfoam\_datapipe.py} (the documented extension point), and
\emph{possibly} the three surface loss functions in \texttt{train.py} if variable
ordering differs (\cref{rem:ordering}). That is a small, honest delta and it should be
legible as such. \texttt{apply.sh} must be idempotent --- you will re-clone --- and
\texttt{UPSTREAM.md} must record the commit the patches were written against, because
they will silently rot when the upstream file changes.
\end{remark}

\begin{remark}[No VTU --- with one bounded exception. A factor of $\sim$1000 in storage]
\label{rem:novtu}
DrivAerML volume meshes run to $\sim$150M elements per case and the full dataset is on
the order of tens of TB. This study is \textbf{surface-only for every metric}: the
surface meshes are $\sim$10M elements, and pressure plus $x$-wall-shear integrated over
the surface is \emph{all} that drag requires (\cref{eq:forces}). Downloading volumes
``in case we want wake analysis later'' is how this project dies in Phase 0.

\textbf{The exception is planned, not opportunistic.} In Phase 5, fetch VTU for exactly
\textbf{five} DrivAer cases --- selected \emph{after} A3 completes, once the error
ranking is known --- purely to evidence the mechanism (\cref{sec:wake}). Five cases is
negligible. Five hundred is the failure this remark exists to prevent. The distinction
is that the five are chosen \emph{because} you now know which cases matter; a
speculative bulk download is chosen because you do not.
\end{remark}

\begin{remark}[\texttt{scaling/} is global, not per-arm --- and that is deliberate]
\label{rem:scalingpath}
This directory layout exists to make the leak in \cref{rem:leak} \emph{awkward to
commit}. \texttt{source\_factors.json} is computed once, over the combined
Ahmed+Windsor training set, and applied \emph{unchanged} to DrivAer at A3 test time.
If scaling factors instead lived inside each arm's output directory --- which is the
natural thing to do, and what \texttt{compute\_statistics.py} will encourage --- then
regenerating them for DrivAer would feel like housekeeping rather than like the
methodological error it is. Making the file global and singular means overwriting it is
a conscious act. \textbf{The A1 baseline gets its own factors, because A1 is trained on
DrivAer and no leak is possible there.} A3 does not.

\textbf{A3a gets its own factors too, and this is not a detail.}
\texttt{ahmed\_only\_factors.json} is computed over Ahmed alone. Reusing
\texttt{source\_factors.json} for A3a would leak Windsor's statistics into the
Ahmed-only arm --- and Windsor is precisely the variable A3a exists to remove. The two
files sit side by side so that grabbing the wrong one is a visible mistake rather than
an invisible one.
\end{remark}

\begin{remark}[\texttt{drivaer\_holdout/} is a separate directory for the same reason]
The recovery curve is only meaningful if evaluated on DrivAer cases that were never
fine-tuned on, at any $N$, under any seed. Keeping the holdout in its own directory,
never referenced by any \texttt{ft\_*} config, makes contamination a thing you would
have to do on purpose. Reserve it in Phase 0 and do not touch it until Phase 5.
\end{remark}

\begin{remark}[\texttt{results/} is append-only JSON, not a notebook]
\label{rem:results}
Four arms plus six fine-tuning budgets times three seeds is 22 runs
(\cref{crit:p4}). Accumulating these in a notebook's memory means losing them to a
kernel restart, and means figures that cannot be regenerated. One JSON per run, matching
the schema in Section 11.4, written immediately on completion; \texttt{figures.py} reads
the directory. A figure is then always reproducible from what is on disk, which is
the difference between a result and an anecdote.
\end{remark}

\subsection{Keeping the plan current --- the memory problem}
\label{sec:memory}

\begin{remark}[This document is the project's memory. Update it.]
Across sessions, and across collaborators, nothing persists except what is written
down. The plan as first written is a \emph{prediction}; the project is a sequence of
\emph{findings}. If, at Phase 4, someone is handed the unmodified Phase-0 document,
they will help you start Phase 0 again.

After each phase: update the Session Restart Card (\cref{sec:restart}); record what was
actually run, what came out, and what broke; and --- this is the part people skip ---
\textbf{adjudicate the hypotheses}.
\Cref{hyp:degrade,hyp:localize,hyp:cancel,hyp:downward,hyp:diversity}
are registered predictions. Mark each held or refuted, with the number that decided it.
That record, more than any figure, is what makes the study honest, because it is what
prevents a prediction from being quietly rewritten after the fact.
\end{remark}

\subsection{Naming caveat: Modulus vs.\ PhysicsNeMo}

The DoMINO paper refers to \textbf{NVIDIA Modulus} and its code links point at
\texttt{github.com/NVIDIA/modulus}. The project was subsequently renamed
\textbf{PhysicsNeMo}. They are the same codebase. Tutorials, forum answers, and paths
containing \texttt{modulus} are the older name, not a different project. Pin your
version and record the commit hash in \texttt{environment.md} --- ``which release?'' is
the first question anyone will ask when your numbers fail to reproduce, and it is the
first question you will ask yourself.

% ============================================================
\section{Session Restart Card}
\label{sec:restart}
% ============================================================

\begin{lstlisting}[style=sh]
PROJECT: Cross-geometry generalization of DoMINO (NVIDIA PhysicsNeMo).

QUESTION: Can DoMINO, trained only on simple bluff bodies (AhmedML +
WindsorML), predict drag on a realistic road car (DrivAerML) it has
never seen? If not, how many DrivAer runs does it take to fix?

NOT the same as the paper's OOD (drag outside training range, SAME body).
Ours is a DIFFERENT BODY. See Remark 5.6.

HEADLINE:   Cd error in drag counts on cold upward transfer.
COMPANIONS: Spearman rho (design ranking); area-weighted Cp error map
            with SIGNED BIAS (error-cancellation detector).

ARMS (run IN ORDER):
      A1  DrivAer -> DrivAer holdout     (instrument check)
      A2  DrivAer -> Ahmed+Windsor       (control; GATE -- run before A3.
                                          If this fails, A3 is
                                          uninterpretable. Stop.)
      A3a Ahmed only -> DrivAer          (interpretive control, H3.5.
                                          OWN scaling factors -- do NOT
                                          reuse A3's. See Rem 15.4)
      A3  Ahmed+Windsor -> DrivAer       (THE experiment)
      + recovery: fine-tune A3 on N in {0,5,10,25,50,100}, 3 seeds each
      = 22 runs total

A3a IS NOT A CURVE. Two points. Plot as bare points on F1, never joined.
The Ahmed(500) vs Ahmed+Windsor(855) sample-count confound is ACCEPTED
and STATED, not closed. Do not claim a diversity effect the design
cannot separate. (Rem 11.1)

PHASES:
  0  Recon + repo structure + compatibility table     [ ]
  1  Baseline reproduction on DrivAerML  (R2 > 0.95)  [ ]
  2  Force integration + metrics harness (<2% error)  [ ]
  3  Unified multi-body datapipe (assertion passes)   [ ]
  4  Four arms + recovery curve (22 runs)             [ ]
  5  Error localization + WAKE EVIDENCE (5 VTU cases) [ ]
     -> adjudicate hypotheses
  6  Writeup: 3 figures, 1 sentence                   [ ]

DEFERRED ON PURPOSE (do not let me scope-creep into these):
  fidelity transfer | UQ | volume-field study | physics losses
  encoder-vs-aggregation ablation (Rem 13.1) -- FOLLOW-UP paper
  diversity SCALING study (how many body families?)  -- FOLLOW-UP paper
    (A3a gives TWO points. It is a control, not a scaling law.)

HYPOTHESIS LEDGER  (registered in advance -- mark held/refuted, with
                    the number that decided it. Do NOT rewrite them.)
  H3.1 upward transfer degrades          [ ] held / refuted:  ______
  H3.2 error localizes to absent features[ ] held / refuted:  ______
  H3.3 error cancellation precedes acc.  [ ] held / refuted:  ______
  H3.4 downward transfer succeeds (ctrl) [ ] held / refuted:  ______
  H3.5 source diversity buys transfer    [ ] held / refuted:  ______
       (A3a cold Cd vs A3 cold Cd)

STATUS: PHASE 0 COMPLETE (2026-07-11). Exit criterion 7.1 cleared.

UPSTREAM COMMITS (pinned, in environment.md):
  physicsnemo         @ 59aaf59b48901bd8df2c9bad83f6d05ea47d8c04
  physicsnemo-curator @ 86533e581b3550326d89e97cb4d4126e7061b416
  physicsnemo-cfd     @ 0d2305e1777351569b1795ce38884ee945491d28

PHASE 0 FOUND SEVEN SILENT TRAPS. None were in the plan. None raise
an error. All produce a converging loss and a worthless model.
Read Rem 7.2 (rem:p0value) before doing ANYTHING else.

  1. loss.py:423 -- channel 0 MUST be pressure, channel 1 MUST be
     x-wall-shear. Hardcoded positional indexing.
  2. NO retraining.py. Fine-tune via resume_dir + train.py -- but it
     restores the EPOCH COUNTER. A 500-epoch checkpoint with
     train.epochs=50 exits immediately having trained NOTHING, and
     writes a checkpoint. RECOVERY CURVE WOULD BE A FLAT LINE.
     Patch #1 = weights-only load + force init_epoch=0.
  3. force_mom_<i>.csv = CONSTANT ref in Ahmed/Windsor,
     VARIABLE ref in DrivAer. Same filename, opposite meaning,
     across exactly the boundary we transfer over. 47 drag counts
     on Ahmed. USE src/datasets.py -- never hardcode the filename.
  4. Windsor surface is .vtu with POINT data. Ahmed/DrivAer are
     .vtp with CELL data. A *.vtp glob finds nothing in Windsor.
  5. Windsor wall shear = 3 separate scalars, file order cfx, cfz,
     cfy -- Z BEFORE Y. Stacking in array order permutes channels.
  6. U_inf spans 39x (Ahmed 1 m/s, DrivAer 38.889). Dimensional
     wall shear spans ~1500x. Must non-dimensionalize or DrivAer
     dominates the loss entirely.
  7. Windsor is Y-UP (not z-up) and YAWED -2.5 deg. Drag is safe
     (x is streamwise in all three). LIFT AND Cf_y/Cf_z ARE NOT
     COMPARABLE. Phase 3 must decide on an axis swap.

DEFUSED: streamwise axis IS x in all three. Verified, not assumed.

SINGLE SOURCE OF TRUTH: src/datasets.py. Every field name, reference
area, U_inf, and file convention lives there, with provenance.
NOTHING ELSE MAY HARDCODE THEM.

NEXT: Phase 1 (baseline reproduction on DrivAerML, Drag R2 > 0.95).
      BLOCKED ON: cluster access. environment.md cluster fields are
      all TODO -- scheduler, GPU count, wall-clock limit, scratch
      quota and purge policy. The 22-run matrix cannot be sized
      without them.

LAST RESULT: recon_all.json (1 case from each dataset, on disk at
             D:\data\{ahmedml,windsorml,drivaerml}\run_1)
BLOCKER: cluster access not yet configured.
\end{lstlisting}

% ============================================================
\section{Collaboration Contract (Paste This First)}
\label{sec:contract}

\begin{remark}[What this is for]
An assistant reading this document does not need help \emph{understanding} it --- it
will read the whole thing. What it needs is a statement of \textbf{how to behave}: what
to verify before acting, what to refuse, and what to push back on. The failure modes
below are not hypothetical; they are the specific ways a well-meaning collaborator
degrades this study. This section is an operating contract, not an orientation.
\end{remark}

\subsection{The prompt}

\begin{lstlisting}[style=sh,basicstyle=\ttfamily\scriptsize]
You are collaborating on the project specified in the attached document.
Read all of it before responding. Then follow this contract.

=== FIRST, ESTABLISH STATE ===
Do not assume we are at Phase 0. Read the Session Restart Card (App. B)
and the Hypothesis Ledger. If STATUS is blank or looks stale relative to
what I'm telling you, ASK before proceeding. Helping me restart a phase
I already finished is a real and costly failure -- you have no memory of
prior sessions, and this document is the only state that persists.

=== DIVISION OF LABOUR ===
I run the compute, the data, and the cluster. You do NOT have GPU access,
dataset access, or a persistent filesystem. You own: code, reasoning,
debugging, and honesty checks. Do not offer to "run" anything. Do ask me
to paste outputs, error traces, JSON, and numbers.

=== THINGS TO REFUSE, EVEN IF I ASK ===
These are not preferences. They void the study.

1. SCALING-FACTOR LEAK (Remark 9.1). surface_factors are computed on the
   SOURCE set ONLY and applied UNCHANGED to DrivAer. If I propose
   recomputing them on DrivAer -- for any reason, however sensible it
   sounds -- refuse and cite Remark 9.1. This is the error most likely to
   survive review and invalidate the headline claim.
   COROLLARY: A3a (Ahmed-only) has its OWN factors, computed on Ahmed
   alone. If I reuse A3's Ahmed+Windsor factors for A3a, Windsor's
   statistics leak into the arm whose entire purpose is to REMOVE
   Windsor. Stop me.

2. HOLDOUT CONTAMINATION. drivaer_holdout/ is never fine-tuned on, at any
   N, under any seed. If a config would touch it, stop me.

3. HYPOTHESIS REWRITING. H3.1-H3.5 (Sec. 6) are REGISTERED PREDICTIONS.
   They were written before the experiment. If results contradict them,
   the hypothesis is REFUTED -- it is not amended, softened, or
   reinterpreted. A refuted hypothesis is a finding. A rewritten one is
   a lie. Hold the line on this even if I want to move it.

4. SKIPPING EXIT CRITERIA. Each phase has one. They exist because the
   dominant failures here are SILENT. If I want to move to Phase N+1
   without clearing Phase N's criterion, say so plainly and make me
   decide consciously. Note A2 is a GATE: if the downward control fails,
   A3 is uninterpretable and there is no point running it.

5. OVERCLAIMING NOVELTY. The DoMINO paper ALREADY tests out-of-
   distribution samples (Remark 5.6) -- but along a different axis (drag
   outside training range, SAME body). Ours is a different BODY. If I
   draft anything implying we invented OOD testing for DoMINO, correct it.

6. OVERCLAIMING FROM A3a. A3a vs A3 is TWO POINTS (Remark 11.1). It says
   whether adding one body family moves the number. It does NOT say how
   many families are needed, and it does NOT separate diversity from
   sample count (Ahmed=500, Ahmed+Windsor=855). If I draft a sentence
   implying a diversity scaling law -- or draw a line through two points
   on F1 -- correct it.

=== WHAT TO BE SUSPICIOUS OF ===
- A cold-transfer (A3, N=0) Cd that looks GOOD. The plan predicts
  degradation. If the number is suspiciously fine, check leakage FIRST
  (Remark 9.1), then normalization, then the force integrator. A neural
  operator does not spontaneously learn what a wheel does to flow.
- A3a and A3 cold numbers that are IDENTICAL to several digits -> suspect
  the same scaling factors were used for both, or the same checkpoint was
  evaluated twice. They are different training runs on different data;
  they should not agree exactly.
- A loss curve that converges beautifully but results that make no sense
  -> suspect variable ordering (Remark 4.1 / 5.5). The loss will look
  healthy while regressing the wrong channels.
- Cd within tolerance but Spearman rho poor -> the model cannot rank
  designs; it is unusable for optimization regardless of MAE (Def. 3.2).
- Large Cp L2 with near-zero signed bias -> error cancellation
  (Remark 3.2). This is H3.3. Report it; do not "fix" it away.

=== TONE ===
Tell me when I am wrong. Tell me when a result smells off. Do not
validate a number just because I sound confident about it. If a claim in
my draft is stronger than the evidence, say so before a reviewer does.

=== WHAT I WANT FROM YOU THIS SESSION ===
[fill in: e.g. "Phase 0. Here is my recon.py JSON output for the three
datasets -- help me fill the compatibility table and pin the variable
ordering." ]
\end{lstlisting}

\subsection{Why each clause exists}

\begin{center}\small
\begin{tabular}{@{}p{3.6cm}p{5.2cm}p{6.4cm}@{}}
\toprule
\textbf{Clause} & \textbf{Failure it prevents} & \textbf{Why it needs stating}
\\
\midrule
Establish state first
  & Restarting a completed phase
  & The assistant has no memory across sessions. This document is the
    \emph{only} persistent state, and a stale paste looks identical to a fresh one. \\[0.3em]
Division of labour
  & Wasted turns; false offers to execute
  & No GPU, no data, no filesystem persistence. Useful on code and reasoning;
    useless on compute. \\[0.3em]
Refuse the scaling leak
  & Voided headline result
  & \cref{rem:leak}. Recomputing on the target is the \emph{natural} thing to do and
    feels like housekeeping. A compliant assistant will happily help. The A3a corollary
    is subtler still: reusing A3's factors silently reintroduces the variable A3a
    exists to remove. \\[0.3em]
Refuse holdout contamination
  & Meaningless recovery curve
  & The curve is only interpretable on cases never fine-tuned on. \\[0.3em]
Refuse hypothesis rewriting
  & Post-hoc rationalization
  & The registered-prediction discipline is the study's main epistemic asset. It is
    also the easiest thing to quietly abandon when results disappoint. \\[0.3em]
Refuse skipping exit criteria
  & Silent corruption propagating
  & Every dominant risk here (\cref{rem:ordering,rem:leak}) fails \emph{quietly}. The
    criteria are the only detectors. A2 in particular is a gate, not a formality. \\[0.3em]
Refuse novelty overclaim
  & Immediate loss of credibility
  & \cref{rem:ood}. Anyone who has read Section 4 of the paper will catch it. \\[0.3em]
Refuse overclaiming from A3a
  & A scaling law asserted from two points
  & \cref{rem:a3a}. A3a is the most \emph{tempting} result in the study to over-read,
    precisely because it looks like the start of a curve. It is not one, and the
    sample-count confound is not closed. \\[0.3em]
Suspect a good cold number
  & Celebrating a leak
  & \cref{hyp:degrade} predicts failure. A good result is more likely a bug than a
    breakthrough --- \emph{check the bug first}. \\
\bottomrule
\end{tabular}
\end{center}

\begin{remark}[The most important clause is the one about suspicion]
An assistant's default failure mode is agreeableness. If a cold-transfer number comes
back looking good and the human sounds pleased, the path of least resistance is to
congratulate and move on. The registered hypotheses exist precisely so that a good
result is \emph{surprising} and therefore \emph{suspect}. Make the assistant hold that
line --- and notice that this cuts both ways: it is also a line \emph{you} will be
tempted to cross, at 2am, when the number finally looks right.
\end{remark}

% ============================================================
\section{Sources}
% ============================================================

\begin{itemize}
  \item Ranade et al., \emph{DoMINO: A Decomposable Multi-scale Iterative Neural
  Operator for Modeling Large Scale Engineering Simulations}, NVIDIA.
  \url{https://arxiv.org/abs/2501.13350}
  \item DoMINO example documentation --- config keys, custom-dataset extension,
  retraining recipe, published validation metrics. \\
  \url{https://docs.nvidia.com/physicsnemo/25.11/physicsnemo/examples/cfd/external_aerodynamics/domino/README.html}
  \item Ashton et al., \emph{DrivAerML}. \url{https://arxiv.org/abs/2408.11969}
  \item Ashton et al., \emph{AhmedML}. \url{https://arxiv.org/abs/2407.20801}
  \item Ashton et al., \emph{WindsorML}. \url{https://arxiv.org/abs/2407.19320}
  \item CAE ML Datasets collection. \url{https://caemldatasets.org}
  \item PhysicsNeMo. \url{https://github.com/NVIDIA/physicsnemo}
  \item PhysicsNeMo-Curator (preprocessing).
  \url{https://github.com/NVIDIA/physicsnemo-curator}
  \item PhysicsNeMo-CFD (force computation, benchmarking).
  \url{https://github.com/NVIDIA/physicsnemo-cfd}
\end{itemize}

\end{document}